%% Before the class: makes Computer Modern scalable, which the watermark needs on a machine where
%% acmart falls back to it because the libertine package is missing. A no-op with Libertine.
\RequirePackage{fix-cm}
%% `review` numbers the lines, for annotated reviews of the PDF (see the REVIEW-*.md files).
\documentclass[sigconf,review]{acmart}

%% EVERY number in this paper - prose, tables, and abstract alike - is a LaTeX macro defined in
%% figures/variables.tex, generated by research/analysis/paper_variables.py. Nothing below should
%% contain a bare numeric literal that came from a measurement: a number appearing in both a table
%% and the prose discussing it used to be typed twice, so a data refresh could update one and miss
%% the other. See that script's module doc comment for which blocks are read back from a
%% measurement artifact and which are transcribed with their provenance recorded alongside them.
%%
%% The empirical-study numbers (Table 1, corpus size, bytes-AST correlation) were measured over
%% the full \NumRepos{}-repository list on 2026-09-07; see research/data/corpus_stats/PROVENANCE.md
%% for the run. They are no longer the 100-repository `small` sample, and this comment said
%% otherwise until 2026-09-09.
\input{figures/variables.tex}

%% Result boxes. Each research question stated in the introduction is answered by a correspondingly
%% numbered research answer, RA1.1--RA5.3, set in one of these at the point where its supporting
%% measurement is reported. A box holds the answer statement only - two or three sentences, every
%% measured number a macro from figures/variables.tex - with the argument behind it following as
%% ordinary prose. Keeping the boxes short is a two-column-layout constraint, not a style
%% preference: a boxed half-page of prose splits badly across a column break.
\usepackage[most]{tcolorbox}
\usepackage{array}
\usepackage{forest}
%% The draft watermark, on every page. Edit or blank \WatermarkText to change or drop it; the
%% package stays so the change is one line. Light and diagonal, so it reads as a stamp, not as a
%% figure, in print and in greyscale. Grey 0.9 rather than real transparency: draftwatermark has
%% no alpha option, and a light grey reads the same on paper and on an e-ink tablet.
\usepackage{draftwatermark}
\newcommand{\WatermarkText}{\sffamily\bfseries PRE-DRAFT}
\DraftwatermarkOptions{text=\WatermarkText, color={[gray]{0.9}}, angle=50}
\newtcolorbox{resultbox}[1]{%
  breakable,
  enhanced,
  colback=black!3,
  colframe=black!55,
  colbacktitle=black!10,
  coltitle=black,
  fonttitle=\bfseries,
  boxrule=0.4pt,
  left=5pt, right=5pt, top=4pt, bottom=4pt,
  title={#1},
}

\setcopyright{acmlicensed}
\copyrightyear{2026}
\acmYear{2026}
\acmDOI{XXXXXXX.XXXXXXX}
\acmConference[Conference'26]{Make sure to enter the correct conference title from your rights
  confirmation email}{June 03--05, 2026}{Ludbreg, Croatia}
\acmISBN{978-1-4503-XXXX-X/26/06}

\begin{CCSXML}
<ccs2012>
   <concept>
       <concept_id>10011007.10011006.10011073</concept_id>
       <concept_desc>Software and its engineering~Software maintenance tools</concept_desc>
       <concept_significance>500</concept_significance>
       </concept>
   <concept>
       <concept_id>10011007.10011074.10011099.10011693</concept_id>
       <concept_desc>Software and its engineering~Empirical software validation</concept_desc>
       <concept_significance>300</concept_significance>
       </concept>
 </ccs2012>
\end{CCSXML}

\ccsdesc[500]{Software and its engineering~Software maintenance tools}
\ccsdesc[300]{Software and its engineering~Empirical software validation}


\keywords{content-aware diffing, edit scripts, diff visualisation, source code differencing, binary file differencing, ground-truth dataset, empirical software engineering}

\title{Content-Aware Diffing: Edit Scripts, Visualisations, and a Living Ground Truth}

\author{Marko Ivankovi\'{c}}
\affiliation{%
  \institution{Universit\"{a}t Passau}
  \city{Passau}
  \country{Germany}
}
\email{marko@ivankovic.me}

\author{Goran Petrovi\'{c}}
\affiliation{%
  \institution{Universit\"{a}t Passau}
  \city{Passau}
  \country{Germany}
}
\email{stgpetrovic@gmail.com}

\author{Ren\'{e} Just}
\affiliation{%
  \institution{University of Washington}
  \city{Seattle}
  \country{USA}
}
\email{rjust@cs.washington.edu}

\author{Gordon Fraser}
\affiliation{%
  \institution{Universit\"{a}t Passau}
  \city{Passau}
  \country{Germany}
}
\email{gordon.fraser@uni-passau.de}

\begin{document}

\begin{abstract}
Diffing, the process of comparing two versions of a file and showing what changed between them,
underlies version control and code review. The files that commits change hold content of many
types: text, code, pictures, archives, documents, fonts. For each type a diffing tool does two
things that are usually treated as one. It computes a \emph{delta}, an edit script by which a
machine converts one version into the other, which is judged by its correctness and its cost. It
also produces a \emph{visualisation}, an account that a human reads in order to understand what
happened, which can be judged only against human reading. We study the distance between the two on
real changes to all content, with code as the content type that has the deepest ground truth.

We built a living ground-truth dataset, the test corpus of an open-source project with
\ProjectStars{} GitHub stars, to which ground truth was added in each of the \GrowthWeeks{} weeks
since its first hand mapping. It holds \NumFixtures{} real code changes in \NumFixtureLanguages{}
languages with a human-authored node mapping, \PresetPainted{} of them with human-painted
renderings, and \ContentAccAllJudged{} changes to pictures, cursors, fonts, archives, message
catalogs and PDFs judged by hand, together with \InvariantCount{} invariants that a visualisation
has to obey and a delta does not. Git judges \ChangeCensusBinaryPct{}\% of the
\ChangeCensusChanges{} files that recent commits of \NumRepos{} repositories change to be binary,
and shows a developer only that they differ. Painted as it stands, the line edit script of
\texttt{diff} breaks at least one invariant on \InvariantLineScriptFixturesPct{}\% of changes, and
the human-preferred tree mapping, painted node by node, still does on
\InvariantTreeScriptFixturesPct{}\%. On \AmbiguityListPct{}\% of code changes the preferred
correspondence cannot be expressed by the common operations (insert, delete, update and move);
\InvisibleNodeShare{}\% of syntax-tree nodes never reach the screen; and
\ContentVerdictNothingPct{}\% of the judged changes to other content, every one a non-empty byte
delta, show a reader nothing that anyone chose to change. A visualisation is also not unique:
\PresetBothPct{}\% of painted changes have a sparse and a generous correct rendering, the generous
one painting \PresetFullOverMinimal{} times the bytes. A whole-tree tree-edit-distance computation
completes within one second for only \RqOneCodePct{}\% of \RqOnePairsAttempted{} sampled code file
pairs, and no configuration of eight established tools renders more than \SrcDiffPerfectPct{}\% of
code changes perfectly.

As a tool contribution, we present \textsc{OmniDiff}, which diffs a file by what it holds: text,
code in \NumSupportedLanguages{} languages, and six families of other content. It diffs
\ContentCensusShownPct{}\% of the changed files as text or by their content. On code, it maps
every visible node as the human did on \MappingZeroVisiblePct{}\% of changes, and on
\ContentAccAllJudged{} judged changes to other content its verdict of how much changed matches a
human's on \ContentAccAllPerfectPct{}\%.
\end{abstract}

\maketitle

\section{Introduction}
\label{sec:intro}

\begin{figure}[t]
  \small
  \begin{tabular}{@{}p{0.45\linewidth}@{\hspace{0.05\linewidth}}p{0.45\linewidth}@{}}
    \textsc{Before} & \textsc{After} \\
    \midrule
    \begin{minipage}[t]{\linewidth}
\begin{verbatim}
def save(rec):
  db.begin()
  db.put(rec)
  db.commit()
\end{verbatim}
    \end{minipage}
    &
    \begin{minipage}[t]{\linewidth}
\begin{verbatim}
def save(rec):
  if rec.dirty:
    db.begin()
    db.put(rec)
    db.commit()
\end{verbatim}
    \end{minipage} \\
  \end{tabular}
  \caption{Wrapping a block in a condition. Only whitespace (indentation) changed, but a
  line-based diff reports all of it as deleted and reinserted. An AST diff can report one
  insertion and a move.}
  \label{fig:motivation}
\end{figure}

Diffing, the process of comparing two versions of a file and showing what changed between them,
underlies modern version control and, especially, modern code review. During a review, developers
read the diff to see what changed and what stayed the same, so the quality of the diff directly
shapes the ease and quality of the review. Whether that reading succeeds depends on how the change
is shown, not only on the change. Reviewers name understanding the change as the main challenge of
code review~\cite{bacchelli-bird-2013}, and ask for tools that show a change's structure rather
than its text delta~\cite{tao-2012}. Controlled experiments show the presentation alone affecting
the outcome. Files shown later in a review receive fewer comments, and a seeded defect in the last
file has 64\% lower odds of being found than the same defect in the
first~\cite{fregnan-2022,rahman-2026}. The unified view of a diff costs less visual effort and
finds more bugs than the split view of the same diff~\cite{alcocer-2023}. A change decomposed
into its separate concerns is reviewed with fewer wrongly reported issues than the same change
tangled~\cite{herzig-2013,barnett-2015,dibiase-2019}. A reviewer's working-memory capacity
predicts whether they find a \emph{delocalised} defect, one whose parts the diff shows far
apart~\cite{baum-2019,goncalves-2022}.

The cost of a diff that is hard to read is paid in time and in quality. At Google, useful
comments fall and review latency rises with the size of a change~\cite{sadowski-2018}, and
Mozilla's contributors name size-related factors as the largest driver of review
latency~\cite{kononenko-2016}. Components with low review participation ship more post-release
defects~\cite{mcintosh-2016}. Practitioners rate the reaction time of review tooling as very
important~\cite{kudrjavets-2023}, and at Meta, developer satisfaction with code review
falls with the time the slowest quarter of a developer's changes wait for
review~\cite{meta-review-time-2022}. Where a diff shows a refactoring as a refactoring rather
than as deleted and inserted lines, reviewers report lower cognitive effort and read fewer
lines~\cite{ge-2017,brito-2021}, and syntactic deltas modelled on how developers edit are judged
more understandable than the cheapest tree edit~\cite{decker-2020}.

What a commit changes is content of many types. In the corpus of Section~\ref{sec:dataset}, the
files that recent commits change are text of every kind (code, data, configuration and
documentation), and they are also pictures, which change in \ContentCensusPicturesRepos{} of
\ContentCensusRepos{} repositories, archives, PDFs, fonts, compiled translation catalogs and
cursor themes. The diff that almost all developers and platforms use today treats every one of
them the same way. A file is a sequence of text lines, and the tool computes an edit script
between the two sequences. A file whose bytes git judges binary, chiefly by a NUL byte, is
reported as \texttt{Binary files differ} and nothing more~\cite{gitattributes}; that is
\ChangeCensusBinaryPct{}\% of the files the commits change, and a reviewer of such a change learns
that it happened and not what it was. For text the line edit script is cheap to compute and
correct, and it is still a poor account of many changes. A whole class of common code
changes—extracting code into a function, wrapping a block in a condition, reordering
declarations—is a move of existing code, which a line script can state only as a deletion plus an
insertion. Figure~\ref{fig:motivation} shows a small example: wrapping three unchanged lines in a
conditional, which a line-based tool reports as three lines deleted and four inserted, and which a
reader has to diff again in their head to see that only the wrapper is new.

We argue that these are two faces of one problem, and that it is not specific to code. An edit
script, which we call a \emph{delta}, is a recipe by which a machine converts one version of a
file into the other. It is judged by whether it is correct and by how little it costs, and a line
script, a tree edit script and a binary delta are all good at exactly that. A \emph{visualisation}
is what a human needs in order to understand what happened, and it is judged against human
reading. The two come apart for every content type. The cheapest tree edit script pairs two
unrelated \texttt{return} keywords; the smallest byte delta between two PNG files says nothing a
viewer can use; a picture re-encoded without loss has a large delta and no visible change. A
visualisation may even be an account that is not minimal and cannot be executed: a move shown as a
move, a rename shown word by word, whitespace suppressed, a picture's change described by how much
of it a viewer would see. Section~\ref{sec:background} develops the distinction, and the rest of
the paper measures it.

Measuring it needs a ground truth that records what humans consider the right account of a
change, for more than one content type, and that can be kept honest as it grows. We built one as
the test corpus of an open-source diffing tool that is under active development and in use, so the
dataset is not a snapshot: it grows by sampling, by hand curation and from the issues users
report, and \InvariantCount{} machine-checked invariants hold its hand-made answers to one
convention.

We examined five groups of research questions. Each group is answered in its own section, in the
order listed here. Wherever a question can be asked of all content we ask it of all content and
report code as one content type among the others; where our data exists only for code, which has
the deepest ground truth, the question says so.

%% The RQ statements below are restated in full in Sections 4-8 next to their answers. Edit both
%% or neither.
\begin{itemize}
  \item \textbf{Research Question Group 1 (Content):} What content do real-world changes touch?
    How large are real-world changes?
  \item \textbf{Research Question Group 2 (Delta versus visualisation):} How often does a
    delta, painted as it stands, break the rules that a visualisation obeys? How much of a delta
    reaches the reader? What percentage of real-world code changes have a human-preferred node
    mapping that the common edit operations (Insert, Delete, Update, Move) cannot express? What
    percentage have a human-preferred mapping whose cost is higher than that of the cheapest
    mapping found?
  \item \textbf{Research Question Group 3 (Uniqueness):} What percentage of real-world code
    changes have multiple optimal node mappings? What percentage have multiple equally correct
    visualisations, and how do those differ?
  \item \textbf{Research Question Group 4 (Exactness):} What percentage of real-world
    source-code changes can a single, whole-tree tree-edit-distance computation complete within
    one second?\footnote{One second is where Nielsen's response-time limits put the boundary
    between an interaction that keeps a user's train of thought and one that visibly interrupts
    it~\cite{nielsen-usability}.}
  \item \textbf{Research Question Group 5 (State of the art):} What do established tools show
    for each type of content? How accurately do current state-of-the-art code diffing tools find
    the human-preferred visualisation? How fast are they?
\end{itemize}

This paper makes the following contributions:

\begin{itemize}
  \item \textbf{Theory:} It separates the delta from the visualisation of a change, defines
    delta, mapping, painting and the projection from a mapping to a painting for content in
    general, and states which properties each has to satisfy
    (Section~\ref{sec:background}).
  \item \textbf{Dataset:} It presents a living ground-truth dataset: \NumFixtures{} real-world
    code changes in \NumFixtureLanguages{} languages, each with a hand-made node mapping,
    \PresetPainted{} of them with hand-painted renderings in up to two presets, and
    \ContentAccAllJudged{} changes to pictures, cursors, fonts, archives, compiled message
    catalogs and PDFs, each judged by hand for how much it changed. The changes were drawn from
    \NumRepos{} open-source repositories by three sampling methods, and a fourth part is the
    Defects4J~\cite{defects4j} changes of the node-mapping oracle of Alikhanifard and
    Tsantalis~\cite{astdiffbenchmark}, solved again by us for comparison with prior work. The
    dataset, the tool used to create it and its entire history are freely available
    (Section~\ref{sec:dataset}).
  \item \textbf{Invariants:} It states \InvariantCount{} invariants that every ground-truth
    answer obeys, checked by machine. They are rules of visualisation that no delta needs, so the
    number of them an account of a change breaks measures its distance from a visualisation
    (Sections~\ref{sec:invariants} and~\ref{sec:gap}).
  \item \textbf{Evaluation:} It reports a census of all \ContentCensusChanges{} files the
    corpus's recent commits change, by what each file holds, an analysis of the shape of
    real-world changes that is independent of any implementation, and a quantified comparison of
    eight established diffing tools against the human ground truth
    (Sections~\ref{sec:content}--\ref{sec:tools}).
  \item \textbf{Tool:} It presents \textsc{OmniDiff}, a tool that diffs text, code and six
    families of other content, each by what it holds, built from the results above
    (Section~\ref{sec:omnidiff}).
\end{itemize}

\section{Background: Deltas and Visualisations}
\label{sec:background}

A diffing tool produces two artifacts that the literature and most tools treat as one: a delta,
which a machine executes, and a visualisation, which a human reads. This section reviews the work
on each, defines both for content in general, and lists the tools the paper measures.

\subsection{Deltas}
\label{sec:deltas}

\textbf{Sequences.} The string-to-string correction problem of Wagner and
Fischer~\cite{wagner-fischer-1974} asks for the cheapest sequence of character insertions,
deletions and changes that turns one string into another, and solves it by dynamic programming.
Unix \texttt{diff} applies the same idea to lines: Hunt and McIlroy~\cite{hunt-mcilroy-1976}
compute a longest common subsequence of the two files' lines and report the rest as deleted and
inserted. Myers' $O(ND)$ algorithm~\cite{myers-1986} finds a shortest such script in time that
scales with the size of the script rather than of the inputs, and underlies most line-based tools
in use today.

\textbf{Block moves.} A shortest insert-and-delete script is not the shortest recipe once copying
is allowed. Tichy~\cite{tichy-1984} extends string correction with block moves and shows that the
resulting scripts are shorter than those based on a longest common subsequence, and
Heckel~\cite{heckel-1978} anchors a comparison on the lines that occur exactly once in both files,
which detects blocks that moved.

\textbf{Trees.} Tai~\cite{tai-1979} generalises string correction to ordered labelled trees, with
node relabelling, deletion and insertion as the operations. Zhang and
Shasha~\cite{zhang-shasha-1989} give the classic dynamic-programming algorithm for this tree edit
distance; it is exact, but its keyroot decomposition repeats work across overlapping subproblems.
RTED~\cite{rted-2011} and its successor APTED~\cite{apted-2015} choose a decomposition strategy
per subproblem, which keeps the $O(n^3)$ worst case and lowers the runtime on real trees
markedly; Bille~\cite{bille-2005} surveys the field. Chawathe et al.~\cite{chawathe-1996} split
the problem in two, a split this paper builds on: first find a good \emph{matching} between the
nodes of the two trees, then derive a minimum-cost edit script of insert, delete, update and
\emph{move} operations that conforms to it. Source-code differencing follows that split.
ChangeDistiller~\cite{changedistiller-2007} adapts the matching to abstract syntax trees (ASTs),
and GumTree~\cite{gumtreediff-2014} matches large isomorphic subtrees top-down and containers
bottom-up. Later work shortens the scripts further by finding more moves~\cite{mtdiff-2016}, by
matching names and partial subtrees~\cite{ijm-2018}, by letting unchanged lines constrain the node
matching~\cite{matsumoto-2019}, or by replacing GumTree's costly recovery phase with a simpler
heuristic~\cite{falleri-2024}. Most of this work measures quality by the length of the edit
script; Fan et al.~\cite{fan-2021} compare three of the algorithms differentially and find that
each produces inaccurate mappings for a sizeable share of file revisions.

\textbf{Bytes.} Delta compression drops the reader altogether. The rsync
algorithm~\cite{tridgell-rsync} sends only the parts of a file the receiver lacks,
VCDIFF~\cite{vcdiff-rfc3284} is a portable format of copy and add instructions,
bsdiff~\cite{percival-bsdiff} finds approximately matching regions of two executables by suffix
sorting, and git stores the objects of a packfile as chains of such deltas~\cite{git-pack-format}.
Hunt et al.~\cite{hunt-vo-tichy-1998} compare delta algorithms on real releases by the two
criteria that matter for a delta, size and speed, and report that a byte-level algorithm produces
smaller deltas than line-based \texttt{diff}.

\textbf{Other content.} Content that is not a sequence of characters has deltas of its own kind.
The delta of two pictures is the set of pixels that differ, the delta of two archives is the list
of members added, removed and changed, and the delta of two fonts is a list of tables or glyphs.
Each is as mechanical as a line script, and each is correct by construction.

\subsection{Visualisations}
\label{sec:visualisations}

\textbf{Text and code.} That the display of a difference is a problem of its own is as old as
syntactic differencing: Yang's program comparison tool~\cite{yang-1991} matches two parse trees
and then pretty-prints both programs side by side with the differing parts highlighted. Later
work moves the account of a change further from the edit script and closer to the reader.
Ldiff~\cite{canfora-ldiff-2009} and LHDiff~\cite{asaduzzaman-lhdiff-2013} tell changed lines from
added and deleted ones and follow lines that moved. srcDiff~\cite{decker-2020} decides by rules
modelled on how developers edit whether two differing elements are one construct changed or one
construct replaced by another, and reports deltas that developers judge more understandable than
the cheapest ones. ClDiff~\cite{cldiff-2018} groups fine-grained differences into statement-level
ones and links related ones, LSdiff~\cite{kim-notkin-2009} summarises systematic changes as rules
with exceptions, and refactoring-aware review tools show a detected
refactoring~\cite{refactoringminer-2022} as a refactoring rather than as deleted and inserted
lines~\cite{ge-2017,brito-2021}. The options of everyday tools are visualisation devices over an
unchanged line script: ignoring whitespace, highlighting changed words inside a line, colouring
moved lines.

\textbf{Pictures.} For images the distinction is established. The mean squared error of two
images is a property of their pixel delta, and Wang et al.~\cite{ssim-2004} show that a measure
of structural similarity agrees with human ratings of image quality better than it does. Yee's
perceptual metric~\cite{yee-2004} decides whether two renderings differ \emph{visibly}, for
regression tests of rendering software, and pixelmatch~\cite{pixelmatch}, written to compare
screenshots in tests, counts a pixel as changed only when its perceptual colour
difference~\cite{kotsarenko-2010} passes a threshold and ignores anti-aliased
pixels~\cite{vysniauskas-2009}.

\textbf{Containers.} The diffoscope tool~\cite{diffoscope} of the Reproducible Builds
project~\cite{lamb-reproducible-2022} recursively unpacks archives, transforms binary formats
into a human-readable form and compares the results. It serves a setting in which the delta
itself is the object of interest, because every byte that differs between two builds matters; a
reviewer of a commit asks a different question, namely what was changed on purpose.

\subsection{Definitions}
\label{sec:definitions}

Let $A$ and $B$ be the two versions of a file. A \emph{content model} reads a version as a set of
\emph{elements}, $E(A)$, each of which owns a part of what the file holds: the lines of a text,
the tokens and nodes of a syntax tree, the pixels of a picture, the members of an archive, the
glyphs of a font, the pages of a document, the messages of a catalog.

\textbf{Delta.} A delta from $A$ to $B$ is a sequence $\delta$ of operations from a fixed set
such that applying it to $A$ yields $B$: $\mathit{apply}(\delta, A) = B$. Correctness is its only
requirement. Among correct deltas, quality is a cost $c(\delta)$, the number of operations or the
size of their encoding, and the best delta is a cheapest one.

\textbf{Mapping.} A mapping is a relation $M \subseteq E(A) \times E(B)$ that says which element
of one version corresponds to which element of the other. An element of $A$ that $M$ leaves out is
deleted and one of $B$ is inserted; a pair whose two elements differ is an update, and a pair
whose position among the other pairs differs is a move. A one-to-one mapping and a delta carry the
same information up to the order of execution~\cite{chawathe-1996}, so a mapping is a delta with
its order forgotten, and every cost model on deltas ranks mappings as well. A mapping that a human
writes down need not be one-to-one: Section~\ref{sec:gap} measures how often it is not.

\textbf{Painting.} A painting is what the reader of a diff sees. For text it is a pair of
labellings $P = (p_A, p_B)$ that give every byte of each version either no label or one of
\emph{delete}, \emph{insert}, \emph{update} and \emph{move}, with \emph{delete} only in $A$ and
\emph{insert} only in $B$. For content that is not text it is a set of regions or members with
such labels together with a \emph{verdict}, a statement of how much a viewer would see of the
change (Section~\ref{sec:ground-truth}). A painting is not executable: it labels both versions and
presupposes both.

\textbf{Projection.} A projection $\pi_\theta$ takes a mapping to a painting. Its parameter
$\theta$ is the reader's preference, for which this paper uses two presets: \emph{Minimal}, which
paints only what changed, and \emph{Full}, which also paints the brackets, separators, leading
whitespace and whole identifiers that a change brings with it. The \emph{literal} projection
paints exactly the bytes owned by each deleted, inserted or updated element and nothing else; it
is what a tool does that colours an edit script as it stands.

\begin{table}
  \caption{A delta and a visualisation of the same change.}
  \label{tab:delta-vs-vis}
  \small
  \begin{tabular}{@{}>{\raggedright\arraybackslash}p{0.19\linewidth}>{\raggedright\arraybackslash}p{0.37\linewidth}>{\raggedright\arraybackslash}p{0.37\linewidth}@{}}
    \toprule
    & Delta & Visualisation \\
    \midrule
    Consumer & A machine, which executes it & A human, who reads it \\ \addlinespace
    Must be & Correct: applying it to one version yields the other & Faithful: what it leaves
      unmarked is unchanged for this reader \\ \addlinespace
    Judged by & Cost: the number of operations or encoded size & Agreement with what a human
      would mark \\ \addlinespace
    Unit & An operation on an element & A marked region of what is displayed \\ \addlinespace
    Executable & Yes & No: it presupposes both versions \\ \addlinespace
    Minimal & By definition & Not necessarily: a move, a whole renamed identifier \\ \addlinespace
    Depends on the reader & No & Yes: \emph{Minimal} and \emph{Full} \\ \addlinespace
    Text & Line edit script & Marked lines and words \\ \addlinespace
    Code & Tree edit script & Marked tokens, moves shown as moves \\ \addlinespace
    Pictures & Differing pixels & Changed regions, and how visible the change is \\ \addlinespace
    Containers & Differing bytes or members & Changed members, each shown by what it holds \\
    \bottomrule
  \end{tabular}
\end{table}

\textbf{What each has to satisfy.} Table~\ref{tab:delta-vs-vis} compares the two. A delta has to
be correct and should be cheap, and nothing else is asked of it: many different cheapest scripts
convert $A$ into $B$, and to the machine that applies them they are interchangeable. A
visualisation answers to a reader, and we state its requirements as four properties.
\emph{Faithfulness}: text that the painting leaves unmarked is unchanged and in place, up to what
the reader's preference declares insignificant, such as indentation under \emph{Minimal}.
\emph{Agreement}: the painting does not contradict the mapping it shows, so a token that the
mapping pairs with an identical token is not painted as deleted on one side and inserted on the
other. \emph{Form}: what is marked follows the units a reader perceives, so a token is marked
whole, a statement takes its terminator, and no highlight ends in trailing whitespace.
\emph{Preference}: the paintings for different readers are ordered and consistent, so
\emph{Full} marks everything \emph{Minimal} marks, and the two never disagree about what
survived. None of the four follows from correctness and minimal cost, and some contradict
minimality: a rename painted as one whole identifier marks bytes that did not change.
Section~\ref{sec:invariants} turns the four properties into \InvariantCount{} checkable
invariants, and Section~\ref{sec:gap} uses them to measure how far a delta is from a
visualisation.

\begin{figure}
  \small
  \setlength{\fboxsep}{0pt}
  \newcommand{\Ins}[1]{\colorbox{black!22}{\strut\texttt{#1}}}
  \newcommand{\Del}[1]{\fbox{\strut\texttt{#1}}}
  \newcommand{\Sp}{\mbox{\ \ }}
  \begin{tabular}{@{}p{0.46\linewidth}@{\hspace{0.06\linewidth}}p{0.46\linewidth}@{}}
    \textsc{Line edit script} & \textsc{Tree edit script} \\
    \midrule
    \begin{minipage}[t]{\linewidth}\ttfamily
      def save(rec):\\
      \Ins{\Sp if rec.dirty:}\\
      \Ins{\Sp\Sp db.begin()}\\
      \Ins{\Sp\Sp db.put(rec)}\\
      \Ins{\Sp\Sp db.commit()}
    \end{minipage}
    &
    \begin{minipage}[t]{\linewidth}\ttfamily
      def save(rec):\\
      \Sp\Ins{if}\ \Ins{rec}\Ins{.}\Ins{dirty}\Ins{:}\\
      \Sp\Sp db.begin()\\
      \Sp\Sp db.put(rec)\\
      \Sp\Sp db.commit()
    \end{minipage} \\
  \end{tabular}

  \vspace{0.8em}
  \begin{tabular}{@{}p{0.46\linewidth}@{\hspace{0.06\linewidth}}p{0.46\linewidth}@{}}
    \textsc{Minimal painting} & \textsc{Full painting} \\
    \midrule
    \begin{minipage}[t]{\linewidth}\ttfamily
      def save(rec):\\
      \Sp\Ins{if rec.dirty:}\\
      \Sp\Sp db.begin()\\
      \Sp\Sp db.put(rec)\\
      \Sp\Sp db.commit()
    \end{minipage}
    &
    \begin{minipage}[t]{\linewidth}\ttfamily
      def save(rec):\\
      \Ins{\Sp if rec.dirty:}\\
      \Sp\Sp db.begin()\\
      \Sp\Sp db.put(rec)\\
      \Sp\Sp db.commit()
    \end{minipage} \\
  \end{tabular}
  \caption{Four accounts of the change of Figure~\ref{fig:motivation}, on its after side; a
  shaded span is marked as inserted. The line script also marks three lines of the before side
  as deleted. The tree script inserts five tokens and moves the block, which costs nothing and
  is not shown.}
  \label{fig:accounts}
\end{figure}

Figure~\ref{fig:accounts} shows the four accounts for the change of Figure~\ref{fig:motivation}.
The line script is correct and marks three statements that did not change. The tree script is
the cheapest correct one and marks five separate tokens, with the space between two of them left
out. The two paintings follow the conventions of the ground truth of Section~\ref{sec:dataset}: the
condition is one insertion, under \emph{Full} together with the indentation of its line, and the
three statements stay unmarked, because being indented further is all that happened to them.

\subsection{Tools}
\label{sec:background-tools}

The tools measured in Section~\ref{sec:tools} are listed here by the content they address.

\textbf{Text.} \textbf{Unix \texttt{diff}} is the gold standard. It treats each file as a
sequence of lines, computes an edit script between the two sequences with Myers' algorithm, and
reports it as hunks of deleted and inserted whole lines. GNU \texttt{diff} caps the algorithm's
cost with a heuristic by default, so on large inputs with many differences the script it reports
can be longer than the shortest one. It has no move operation and no sub-line detail.
\textbf{\texttt{git diff}} defaults to the \texttt{xdiff} engine, which offers four algorithms:
Myers, the default; \emph{minimal}, which spends more time to find a shortest script;
\emph{patience}, which anchors on lines unique to both sides before recursing; and
\emph{histogram}, a faster generalisation of patience. All four emit the same output form as Unix
\texttt{diff}. \textbf{Neovim's diff mode} (\texttt{nvim -d}) runs the same \texttt{xdiff} line
pass and then a second, character-level pass that highlights what changed inside each paired
line. \textbf{BDiff}~\cite{bdiff} reasons about blocks of lines rather than single lines and
reports character ranges within changed lines.

\textbf{Code.} \textbf{GumTree}~\cite{gumtreediff-2014} is an AST diff. Instead of one global
tree-edit-distance computation, it runs a top-down phase that greedily matches large isomorphic
subtrees, then a bottom-up phase that matches remaining container nodes by a Dice coefficient of
their already-matched descendants, and a recovery pass that pairs up deleted and inserted
identical subtrees the first two phases missed. \textbf{difftastic}~\cite{difftastic} and
\textbf{diffsitter}~\cite{diffsitter} are a newer generation of tools built on
tree-sitter~\cite{tree-sitter} grammars. Both prioritise a diff a human reads comfortably over an
exact node-to-node mapping: difftastic reduces the parse tree to an s-expression and aligns it
with its own structural algorithm, and diffsitter applies a related alignment over a smaller,
hand-picked set of compiled-in grammars. \textbf{srcDiff}~\cite{decker-2020} diffs the srcML
markup of C, C++, C\#, Java and Python source, detects moves, and writes one merged document in
which the text of each side is marked as deleted, inserted or common.

\textbf{Other content.} Git judges a file binary by its bytes, chiefly by a NUL byte, and for a
binary file prints that the two versions differ and nothing more. A user can configure a
\texttt{textconv} driver per path pattern, a program that converts each side to text before the
line diff runs. The conversion is lossy, and each format needs its own program and its own
configuration~\cite{gitattributes}. diffoscope and pixelmatch, introduced in
Section~\ref{sec:visualisations}, are the two tools of this kind that \textsc{OmniDiff} builds
on.

\section{A Living Ground-Truth Dataset}
\label{sec:dataset}

This section describes the dataset the research questions are answered on: the project it belongs
to, the corpus of repositories it is drawn from, how changes were sampled, what ground truth each
content type carries, the invariants that hold the ground truth together, and how the dataset
grows. No research question is answered here.

\subsection{The project and the corpus}
\label{sec:project}

The dataset is the test corpus of \textsc{OmniDiff}, the tool of Section~\ref{sec:omnidiff},
which is developed in the open. On \ProjectStatsDate{} its repository, created in
\ProjectCreated{}, had \ProjectStars{} GitHub stars and \ProjectForks{} forks. Developers star a
repository to show appreciation and to bookmark it~\cite{borges-stars-2018}, so we report the
count as a signal of interest, not of use. The tool is
released as source on crates.io, as prebuilt binaries for Linux, macOS and Windows, through
Homebrew, a signed apt repository and a Nix flake, and with editor integrations for VS~Code and
Neovim. Every fixture of the dataset is also a regression test of the tool, and a release is gated on
the tool's accuracy over all of them. This is what makes the dataset a
living one. It is not frozen for a paper and then left; it is the thing the project's daily
development is measured against, and Section~\ref{sec:growth} shows how it has grown.

To understand the environment diffing tools are used in, we built a corpus of \NumFiles{} files
in \NumLanguages{} languages from two sets of open-source repositories:

\begin{itemize}
  \item The Curated list consists of 100 open-source repositories, selected manually from several
    lists of the most popular repositories and extended by repositories we were already aware of.
    This biases the list towards active, widely used codebases, with good coverage of many
    programming languages. The Curated list is also the smaller list the project's automated
    test infrastructure runs over.
  \item The Full list consists of \NumRepos{} repositories, created by expanding the Curated
    list with all repositories present in the Gentoo Linux distribution's package list. It
    includes codebases of widely different activity, code quality, user base and mix of
    languages.
\end{itemize}

We cloned each repository to a depth of \CorpusCloneDepth{} commits from every main branch. While
we cannot guarantee that all repositories remain accessible, the list of repositories, the cloning
script and all analysis are available from the paper's repository and can be replicated.

We fetched, parsed and measured the corpus on a single machine: an \MachineCpu{} with
\MachineRamGb{}\,GB of RAM, running \MachineOs{}. The checkouts and the measurement database sit
on \MachineDiskCount{} \MachineDiskSizeTb{}\,TB \MachineDiskModel{} drives in \MachineDiskArray{}.
We record this because some results are hardware-bound, notably the wall-clock budget we hold
whole-tree APTED to in Section~\ref{sec:speed} and the per-tool timing of
Section~\ref{sec:tools}.

\subsection{What the repositories hold}
\label{sec:files}

We categorised the files by extension and content. Only two thirds of all files are code
(Figure~\ref{fig:tips}). The rest is data, configuration, documentation, and a residual of
unknown type, which limited manual inspection suggests is mostly data. A tool that diffs only
code leaves a third of the files to something else.

Restricted to code files, the top row of Figure~\ref{fig:corpus-shape} shows the distribution of
bytes, lines of code and AST nodes per file. The median code file has \LocPFifty{} lines, the
99th percentile \LocPNinetyNine{} and the largest \LocMax{}. The node counts are over the
\ParsedCodeFilesPct{}\% of code files that can be parsed (\ParsedCodeFiles{} of \CodeFiles{}); the
rest are in a language without a grammar, flagged as generated and not parsed, or empty. Of those
that parse, the median has \AstPFifty{} AST nodes, the 99th percentile \AstPNinetyNine{} and the
largest \AstMax{}. Byte size and AST node count are strongly correlated for files within the 99th
percentile of both, Pearson $r = \CorrelationRTrimmed{}$, which simplifies to $|\text{AST}|
\approx \text{bytes} / 5$ for quick estimates.

\begin{figure}
  \centering
  \includegraphics[width=0.95\linewidth]{figures/tips}
  \caption{Distribution of file types across the corpus.}
  \label{fig:tips}
\end{figure}

\begin{figure*}
  \centering
  \includegraphics[width=\linewidth]{figures/corpus_shape}
  \caption{The shape of the corpus, every distribution drawn the same way: the share of the
  population at or below the value on the axis, with the median, 90th and 99th percentiles marked
  on each curve and the maximum in the corner. Top: size of every code file in the corpus, in
  lines and bytes, and in AST nodes for the code files that parse, which is why the third
  panel's $n$ is smaller. Bottom: size of every code-file modification in the
  corpus's recent history, in lines changed per file, lines changed per commit, and the share of
  the file an edit rewrites. The horizontal axes are logarithmic except the last.}
  \label{fig:corpus-shape}
\end{figure*}

\subsection{Sampling changes}
\label{sec:diff-shape}

\textbf{Code.} From both the Curated and the Full list of repositories, we randomly sampled
\SampleDiffsPerLanguage{} diffs per language. That gave \SampleCuratedSampled{} diffs from the
Curated list across \SampleCuratedLanguages{} languages, and \SampleFullSampled{} from the Full
list across \SampleFullLanguages{} languages. The two language counts differ because the Curated
list's repositories hold almost no R or Scala code, so those two languages are drawn from the Full
list only. The Full list's count exceeds \SampleDiffsPerLanguage{} per language because some
languages were drawn a second time after their first rounds were largely rejected. Every sampled
diff was manually inspected for sampling errors. At this stage, \SampleRejectedTotal{} diffs were
rejected, mostly because language detection detected the wrong language: in particular TypeScript
(\SampleRejectedTypeScript{} rejected) and R (\SampleRejectedR{}) file extensions were used for
other languages or data. This resulted in a dataset of \SampleCuratedPromoted{} diffs from the
Curated and \SampleFullPromoted{} diffs from the Full list. As Section~\ref{sec:content} shows, a
pure random sample is heavily slanted towards small changes. To account for this, the sample was
augmented by a stratified sample based on the lines of code of the files being changed. The
strata were 1 to 10 lines, 11 to 30, 31 to 100, 101 to 300, 301 to 1000, 1001 to 3000 and more
than 3000 lines of code, and the larger of the two sizes was used. In each stratum,
\SampleStratifiedPerCell{} diffs were randomly sampled per language, and inspected as
above.\footnote{\label{fn:in-progress}The dataset and this paper are being written in parallel.
Drawing and solving the stratified sample, and hand-annotating the Defects4J list below, are both
still in progress, and the ambiguity pass over those two is not finished. Every figure in this
paper is over the \NumFixtures{} changes solved so far, except the cost comparison of
Section~\ref{sec:gap}, the tool comparison of Section~\ref{sec:tools} and the line-level, speed
and robustness figures of Section~\ref{sec:omnidiff}, which were measured at an earlier stage of
the dataset, over the \UnixDiffFixtures{} changes solved then, and, for \textsc{OmniDiff}'s own
rows in them, with an earlier version of the tool. Every figure that depends on the ambiguity
pass is a lower bound for those two datasets.}

A fourth list extends the dataset. The AST node-mapping oracle of Alikhanifard and
Tsantalis~\cite{astdiffbenchmark} contains the \OracleUnits{} AST mappings for \OracleCases{}
bug-fixing changes of Defects4J~\cite{defects4j}. The file contents its offsets index into come
from Falleri and Martinez' replication package~\cite{gumtreesimple}. We solve those same
\OracleUnits{} diffs, so that our mappings can be compared against an independently authored
reference on a corpus that other AST-differencing evaluations already
use.\textsuperscript{\ref{fn:in-progress}}

\textbf{Other content.} From the Full list we sampled changed pairs per content family:
pictures, cursors, fonts, archives, compiled message catalogs and PDFs. Diffing several files at
once and detecting moves across files are out of scope of this paper.

\subsection{Ground truth}
\label{sec:ground-truth}

The ground truth of a change depends on what the file holds. Code has the deepest one, with two
layers that correspond to the mapping and the painting of Section~\ref{sec:definitions}; other
content has a verdict. The annotators are the authors, who each have at least 15 years of
experience in software development and have reviewed tens of thousands of real-world code changes
in industrial environments.

\textbf{Mappings.} Each code change was solved by hand. Solving a change means writing down, for
every node of one AST, which node of the other AST it corresponds to, or that it has none. Where
the annotator judged several nodes of one side interchangeable with several of the other, the
correspondence is recorded as a \emph{multi-map group}, a set of nodes on each side of which any
consistent pairing is equally correct.

\textbf{Paintings.} For \PresetPainted{} of the changes the annotator additionally painted the
rendering: the spans of text a reader should see marked as deleted, inserted, updated or moved,
chosen independently of which nodes carry them. A painted change is one of three cases. In the
first, it carries a \emph{Minimal} and a \emph{Full} painting, the annotator's finding that the
change has a sparse and a generous correct rendering; a tool's sparse output is compared with the
first and its generous output with the second. In the second, it carries a single painting that
names no preset, the finding that the change has one right rendering, which every preset of a
tool has to produce. In the third, it carries a single painting named for one preset, and only
that preset is compared. Section~\ref{sec:uniqueness} counts the cases.

\textbf{Verdicts.} Content other than code has no AST to map, so its ground truth answers a
different question: how much did the change change? The annotators judged \ContentAccAllJudged{}
of the sampled pairs on a six-level scale. Three levels are strengths of ``nothing'':
\emph{invisible}, where only bytes differ (metadata, a lossless re-encoding);
\emph{imperceptible}, visible only once a tool highlights it; and \emph{artifacts}, visible but
chosen by no one, as compression or resampling leaves it. Three are kinds of ``something'':
\emph{edited}, where a part changed deliberately and the rest stayed the same; \emph{redrawn},
where the whole changed and is still a version of the same thing; and \emph{replaced}, where the
after side is not a version of the before side. An archive, a font, a PDF, a catalog or a cursor
of several pictures is judged as a whole and per changed member, \ContentAccVerdicts{} verdicts in
all.

\textbf{Embedded code.} Code also changes inside files of other languages. \EmbeddedFixtures{} of
the \NumFixtures{} code changes edit code held in a string of the host language, in
\EmbeddedLanguages{} embedded languages, and their ground truth maps and paints the embedded
code's own AST as well as the host's.

\subsection{Invariants}
\label{sec:invariants}

A hand-made ground truth of this size is made by several people over months, and its answers
drift: the same kind of change is painted one way in one fixture and another way in the next. We
therefore state the conventions of the ground truth as \emph{invariants}, properties that the
mapping and the paintings of every fixture have to satisfy, each checked by a program that reads
the ground truth alone and runs as a test of every fixture. There are \InvariantCount{} of them.
They were added one at a time, as annotating and reviewing real changes turned up a contradiction
inside a fixture or two fixtures answering one question differently, and each was decided by the
annotators, not derived from a tool's output.

../../../plots/invariants_table.tex

Table~\ref{tab:invariants} shows the invariants in five families, which make the four properties
of Section~\ref{sec:definitions} concrete. Three rules concern the \emph{form} of any painting:
no highlight ends in a row's trailing whitespace, no two painted ranges claim the same byte, and
a token such as \texttt{<=} or \texttt{private} is painted whole. Nine state what each
\emph{preset} paints and how the two relate: \emph{Full} covers everything \emph{Minimal} covers;
a line that is inserted or deleted whole is painted whole under \emph{Full}; \emph{Minimal} never
paints a line's indentation; a renamed identifier is painted at its differing words under
\emph{Minimal} and entire under \emph{Full}; the space before an insertion inside a line goes
with the insertion exactly when the insertion changed how many spaces are there; and the two
presets never disagree about what survived. Ten require the painting to \emph{agree} with the
mapping: a leaf that the mapping pairs with an identical leaf is never painted deleted on one
side and inserted on the other, a leaf the mapping deletes has at least one painted byte, and a
byte painted as moved is never one the mapping leaves unmatched. Four keep \emph{syntactic units}
together: a \texttt{;} goes with its statement in the mapping and in the painting, a bracket with
its partner, and the dot of a call with the member called. Five concern the \emph{form of a
mapping}, for example that two subtrees recorded as identical carry the same tokens. All but the
\InvariantMappingOnlyCount{} rules that read the mapping alone are rules about a painting, and
none is needed for an edit script to be correct. We return to this in Section~\ref{sec:gap}.

The invariants also repair the ground truth. When the set reached \InvariantCount{} rules, the
ground truth of all fixtures, including the hand-written regression fixtures that this paper
does not report on, broke them \InvariantRepairBeforeViolations{} times in
\InvariantRepairBeforeFixtures{} fixtures. We repaired \InvariantRepairedFixtures{} fixtures, each
change being queued for review by an annotator, which left \InvariantRepairAfterViolations{}
violations in \InvariantRepairAfterFixtures{} fixtures. Over the \InvariantHumanAllScored{}
changes this paper reports on, \InvariantHumanAllFixtures{} still carry a violation
(\InvariantHumanAllViolations{} in all), each recorded in the fixture as a known one, and
\InvariantHumanAllCleanPct{}\% carry none. The remainder includes cases in which two rules ask for
different things of the same bytes and cases that await a decision of the annotators.

\subsection{Growth}
\label{sec:growth}

\begin{figure}
  \centering
  \includegraphics[width=\linewidth]{figures/dataset_growth}
  \caption{Growth of the ground truth, by the date each change's ground truth was first
  committed to the project's repository.}
  \label{fig:growth}
\end{figure}

../../../plots/dataset_growth_table.tex

The dataset grows in three ways. Sampling rounds over the corpus are drawn, inspected and solved,
as described above. Hand-written fixtures pin one change pattern each; they are regression tests
rather than a sample, and no figure in this paper includes them. Third, a misbehaviour that a
user reports becomes a fixture with a hand-made ground truth, so that the dataset accumulates the
cases that users of the tool find hard.

Figure~\ref{fig:growth} and Table~\ref{tab:growth} show the growth, read from the history of the
project's repository. The first hand mapping in the present format was committed in
\GrowthFirstMonth{}. In the \GrowthWeeks{} weeks since, ground truth was added or painted in
\GrowthActiveWeeks{} of them, at a median of \GrowthMedianPerWeek{} new changes per week and at
most \GrowthMaxPerWeek{} in one week, to \GrowthCodeFixtures{} mapped code changes,
\GrowthPaintedFixtures{} painted ones and \GrowthContentFixtures{} judged changes to other
content. The dates are those of the first commit of a fixture under its current name, so a
renamed fixture is dated late.

The third way is the newest and the smallest. Of the \GrowthIssuesTotal{} issues in the project's
tracker on \ProjectStatsDate{}, \GrowthIssuesExternal{} were filed by people other than the
authors; most of those ask for features of the viewer or for packaging, and two ask for a further
content type. So far
\GrowthIssuesWithFixtures{} of them has led to ground truth: a request for structural diffing of
SQL led to SQL support in the tool and to \GrowthIssueFixtures{} SQL changes, sampled from the
corpus and solved by hand, which is the last row of Table~\ref{tab:growth} and the dataset's
\NumFixtureLanguages{}th language. The project's mapping viewer also lets a reader file a disagreement with one
node of one fixture's mapping as an issue against that fixture, and two such issues are open. We report the count as it is: the dataset is open to growth from reports, and
almost all of its growth so far is from sampling.

\subsection{Cross-validation with existing datasets}
\label{sec:oracle}

The AST node-mapping oracle of Alikhanifard and Tsantalis~\cite{astdiffbenchmark} records, for
\OracleCases{} bug-fixing changes of Defects4J over \OracleUnits{} Java compilation units, which
before-side node each after-side node corresponds to: \OracleRecords{} node correspondences in
all, built by running six tools, inspecting every diff by hand against six stated criteria, and
resolving disagreements by vote. It was authored by different people, for a different tool, on a
different AST, before \textsc{OmniDiff} existed.

Their oracle uses Eclipse JDT nodes and ours tree-sitter nodes, so the two can be compared only
where a record's character span is one some tree-sitter node also has. \OracleResolutionRate{}\%
of records resolve; the remainder are JDT constructs tree-sitter does not model separately,
chiefly the elements inside a Javadoc comment, which tree-sitter reads as a single token.

The \AmbiguityDefectsFourJTotal{} compilation units we have solved so far are reported as a fourth
dataset, Defects4J, beside the Curated, Full and Stratified samples throughout, and are included
in every pooled total. Over the \OracleHumanUnits{} of them for which the comparison has been
run, our mapping and the oracle's agree to \OracleHumanPrecision{}\%: \OracleHumanDisagreements{}
disagreeing node pairs out of \OracleHumanMappings{}, and \OracleHumanStatementDisagreements{} of
\OracleHumanStatementMappings{} at statement level.

\section{Content of Real-World Changes}
\label{sec:content}

This section measures what a diffing tool is asked to show: which content commits change, and how
large the changes are.

\textbf{RQ1.1.} What content do real-world changes touch?

Over the \CorpusCloneDepth{} most recent commits of each repository of the Full list, we counted
every changed file, binary ones included: \ChangeCensusChanges{} changes. Git judges
\ChangeCensusBinaryPct{}\% of them (\ChangeCensusBinary{}) binary. We then classified by its
bytes each of the \ContentCensusChanges{} changed files that the commits' trees still hold,
leaving out submodule entries; Section~\ref{sec:omnidiff-content} describes the recognition.

\begin{table}
  \caption{Changed files over the corpus's recent history, by what they hold. \emph{Repos.} is
  the number of the \ContentCensusRepos{} repositories with at least one such change.}
  \label{tab:content-census}
  \small
  \begin{tabular}{@{}l>{\raggedright\arraybackslash}p{0.43\linewidth}rr@{}}
    \toprule
    Content & Formats & Changes & Repos. \\
    \midrule
    Text & code, data, documentation & \ContentCensusText{} & -- \\ \addlinespace
    Pictures & PNG, JPEG, GIF, WebP, BMP, ICO, TIFF & \ContentCensusPictures{} & \ContentCensusPicturesRepos{} \\ \addlinespace
    Archives & zip (jar, Office, EPUB), tar, gzip, xz, bzip2 & \ContentCensusArchives{} & \ContentCensusArchivesRepos{} \\ \addlinespace
    PDFs & PDF & \ContentCensusPdfs{} & \ContentCensusPdfsRepos{} \\ \addlinespace
    Fonts & TrueType, OpenType, WOFF, WOFF2, EOT & \ContentCensusFonts{} & \ContentCensusFontsRepos{} \\ \addlinespace
    Catalogs & gettext \texttt{.mo}, Qt \texttt{.qm} & \ContentCensusCatalogs{} & \ContentCensusCatalogsRepos{} \\ \addlinespace
    Cursors & Windows \texttt{.cur}, \texttt{.ani}, X cursors, Hyprland \texttt{.hlc} & \ContentCensusCursors{} & \ContentCensusCursorsRepos{} \\
    \midrule
    Other & binary, or failing to decode & \ContentCensusLeft{} & -- \\
    \bottomrule
  \end{tabular}
\end{table}

Table~\ref{tab:content-census} shows the result. The change counts are decided by a few
repositories: fonts come from \ContentCensusFontsRepos{} repositories and cursors from
\ContentCensusCursorsRepos{}, so the repository count is the one that says how common a format
is. By it, pictures are the most widespread content other than text by far, changed in
\ContentCensusPicturesRepos{} repositories, followed by archives and PDFs. The \emph{Other} row,
\ContentCensusLeftPct{}\% of changes, is mostly text in legacy encodings, compiled objects, audio,
and EOT fonts compressed with MicroType Express. Within text, code is one kind among several:
\EditsCodeSharePct{}\% of the files the commits touched were code files, and the rest of the text
is data, configuration and documentation.

\begin{resultbox}{RA1.1 (Content)}
  \ContentCensusTextPct{}\% of the \ContentCensusChanges{} changed files hold text, of which code
  is one kind, and \ContentCensusFamiliesPct{}\% hold pictures, archives, PDFs, fonts, message
  catalogs or cursors. Pictures change in \ContentCensusPicturesRepos{} of the
  \ContentCensusRepos{} repositories.
\end{resultbox}

\textbf{RQ1.2.} How large are real-world changes?
\label{sec:edit-shape}

We have size data for one content type, code, where size can be counted in lines. Across both
repository lists we analysed \EditsCommits{} commits, at most \CorpusCloneDepth{} commits deep
per repository, so that large repositories, notably the Linux kernel, do not skew the results.
They hold \EditsCodeFileEdits{} code-file edits, of which \EditsModifiedSharePct{}\% were
modifications; the remainder created or deleted a file and so present no pair of versions to
compare. To allow for comparison with existing research, \texttt{git diff} was used to count
changed lines.

The bottom row of Figure~\ref{fig:corpus-shape} shows the sizes of the edits. The median file
edit changes \EditsLinesPerFilePFifty{} lines and the median commit
\EditsLinesPerCommitPFifty{}; \EditsFileEditsUnderTenPct{}\% of file edits and
\EditsCommitsUnderTenPct{}\% of commits change ten lines or fewer, while the 99th percentiles are
\EditsLinesPerFilePNinetyNine{} and \EditsLinesPerCommitPNinetyNine{} lines. Our results are
broadly aligned with existing research, among them Arafat and Riehle, who report that 47.5\% of
all commits touch 10 lines or fewer~\cite{commit-size-distribution-paper}, with the long tail
stretching to tens of thousands of lines. The median edit rewrites only \EditsChurnPFiftyPct{}\%
of the file it touches, and \EditsChurnUnderFivePct{}\% of edits rewrite under 5\%. In the 99th
percentile, an edit rewrites the file entirely (\EditsChurnPNinetyNinePct{}\%).

\begin{resultbox}{RA1.2 (Size)}
  The median code-file edit changes \EditsLinesPerFilePFifty{} lines and rewrites
  \EditsChurnPFiftyPct{}\% of its file; \EditsFileEditsUnderTenPct{}\% of edits change ten lines
  or fewer, and at the 99th percentile an edit rewrites its file entirely.
\end{resultbox}

A diffing tool therefore has to be good at two unrelated jobs: finding a handful of changed lines
in an otherwise untouched file, which is the overwhelming majority of edits, and staying usable
when nothing is left to anchor on. We did not measure the size of changes to other content, for
which a line count has no meaning; the verdicts of Section~\ref{sec:gap} are the closest
measure we have.

\section{Delta versus Visualisation}
\label{sec:gap}

Section~\ref{sec:background} argues that a delta and a visualisation answer to different
requirements. This section measures the distance between them in four ways: by the invariants a
painted delta breaks, by how much of a delta reaches the reader, by the mappings no delta can
express, and by the mappings a cost model ranks below the cheapest.

\textbf{RQ2.1.} How often does a delta, painted as it stands, break the rules that a
visualisation obeys?

The invariants of Section~\ref{sec:invariants} are rules of visualisation, so we hold three
accounts of each change against them. The first is the hand painting itself. The second is the
\emph{line script}: the edit script of GNU \texttt{diff}, with every deleted line painted as
deleted and every inserted line as inserted, from its first column to its last character. The
third is the \emph{tree script}: the human node mapping read as a tree edit script and painted by
the literal projection of Section~\ref{sec:definitions}, that is, every deleted or inserted leaf
painted whole and every leaf paired with a differently reading leaf painted as an update. A move
is free in the cost model and the script does not mention it. The tree script is the most
favourable delta we can construct, because its mapping is the human's own and so has no mapping
error. A delta has no presets, so each script is filed as the \emph{Minimal} and as the
\emph{Full} painting alike. The population is the \InvariantFixtures{} painted changes whose two
files parse and that hold no embedded code, and each script is judged by the
\InvariantPaintingCount{} invariants that read a painting.

The right-hand columns of Table~\ref{tab:invariants} show the result. The line script breaks an
invariant on \InvariantLineScriptFixturesPct{}\% of the changes
(\InvariantLineScriptFixtures{}), \InvariantLineScriptViolations{} times in all. Most of its
violations come from one of two sources. On \InvariantLineScriptSurvivorFixturesPct{}\% of the
changes it marks as deleted and inserted, or as removed, a token that the human mapping keeps,
which is the cost of a script without updates or moves. Most of the rest is indentation, which a
line script always marks and \emph{Minimal} never does. The tree script breaks an invariant on
\InvariantTreeScriptFixturesPct{}\% of the changes (\InvariantTreeScriptFixtures{}),
\InvariantTreeScriptViolations{} times. It almost never contradicts its own mapping: it breaks an
agreement rule on \InvariantTreeScriptAgreementFixturesPct{}\% of the changes. What it lacks is
everything between the nodes: on \InvariantTreeScriptWhitespaceFixturesPct{}\% of the changes it
breaks one of the four rules about which whitespace a painting covers, and most of the remainder are
rules of granularity, such as a renamed identifier painted entire under both presets. The hand
paintings break a painting rule on \InvariantHumanFixturesPct{}\% of the changes.

\begin{resultbox}{RA2.1 (Invariants)}
  Painted as it stands, a line edit script breaks at least one of the \InvariantPaintingCount{}
  painting invariants on \InvariantLineScriptFixturesPct{}\% of \InvariantFixtures{} changes, and
  a tree edit script built from the human-preferred mapping on
  \InvariantTreeScriptFixturesPct{}\%, against \InvariantHumanFixturesPct{}\% for the hand
  paintings.
\end{resultbox}

The two scripts fail differently, and the difference is the argument of this paper in small. The
line script has the wrong correspondence, and no rendering can repair that. The tree script has
the right correspondence and is still not a visualisation on
\InvariantTreeScriptFixturesPct{}\% of the changes, because the
correspondence is over nodes and the reader sees text: which whitespace, which separator and how
much of a changed word to mark is decided by a projection, not by the mapping.

\textbf{RQ2.2.} How much of a delta reaches the reader?

For code, we count the nodes that can be displayed at all. A syntax tree contains two kinds of
nodes. Some carry text of their own, such as an identifier, an operator, a literal or a comment.
Others are pure structure: a \texttt{block} whose extent is exactly its children's, an
\texttt{expression\_statement} wrapping one expression, the chain of single-child nodes a grammar
interposes between a statement and the name inside it. A rendered diff highlights regions of
source text, so a structural node does not appear in the output on its own. We call a node
visible if any byte inside its extent lies outside every one of its children. A tool that maps
every displayed token correctly and every structural wrapper wrongly produces a diff that is
correct, and it scores badly on a metric that counts all nodes equally.

For other content, we count the changes that have nothing to show. Every judged pair of
Section~\ref{sec:ground-truth} differs in its bytes, so each has a non-empty byte delta. Of the
\ContentVerdictPairs{} pairs, the annotators judged \ContentVerdictInvisible{} invisible,
\ContentVerdictImperceptible{} imperceptible and \ContentVerdictArtifacts{} to hold only
artifacts, \ContentVerdictNothing{} in all, and of the \ContentVerdictMembers{} changed members
of containers that they judged separately, \ContentVerdictMembersNothing{}
(\ContentVerdictMembersNothingPct{}\%) fall in the same three levels. The remaining pairs are
\ContentVerdictEdited{} edited, \ContentVerdictRedrawn{} redrawn and \ContentVerdictReplaced{}
replaced.

\begin{resultbox}{RA2.2 (What reaches the reader)}
  \InvisibleNodeShare{}\% of the AST nodes in \NumFixtures{} code changes carry no text of their
  own and cannot appear in a rendered diff. Of \ContentVerdictPairs{} judged changes to other
  content, all of which differ in their bytes, \ContentVerdictNothingPct{}\% show a viewer no
  deliberate change.
\end{resultbox}

The consequence for code is not that node-level metrics should be abandoned. It is that a rate
over all nodes and a rate over visible nodes answer different questions, and any paper reporting
one should state which. This paper reports both. The consequence for other content is that a
tool which reports every pair with a non-empty delta as changed misreports
\ContentVerdictNothingPct{}\% of them.

\textbf{RQ2.3.} What percentage of real-world code changes have a human-preferred node mapping
that the common edit operations (Insert, Delete, Update, Move) cannot express?

The human solvers of Section~\ref{sec:ground-truth} were not able to solve all changes with a
one-to-one mapping. \AmbiguityListWith{} of \AmbiguityListScored{} solved changes
(\AmbiguityListPct{}\%) require an $N{:}M$ mapping: one or more nodes from one side map to one or
more nodes from the other.\footnote{The most common cause of an $N{:}M$ mapping is a refactoring,
where two or more similar blocks of code are moved into a common function. The preferred solution
is the one that maps all of them as moved.} That rate is over the \AmbiguityListScored{} changes
of the Curated and Full lists. The two lists agree within sampling noise, at
\AmbiguityCuratedPct{}\% of \AmbiguityCuratedTotal{} changes on the Curated list and
\AmbiguityFullPct{}\% of \AmbiguityFullTotal{} on the Full list
(Table~\ref{tab:expressibility}). Those two are the datasets whose ambiguity pass is finished, and
the pooled rate RA2.3 states is over them. The table also carries the Stratified sample
(\AmbiguityStratifiedPct{}\%) and the Defects4J changes (\AmbiguityDefectsFourJPct{}\%), and a
rate over all \AmbiguityScored{} changes, but all three are lower bounds: those two passes are
still running, and a change nobody has reviewed for ambiguity is recorded as having none. The gap
between \AmbiguityListPct{}\% and \AmbiguityAnyPct{}\% is the distance annotation still has to
travel, not a difference between the datasets.

\begin{table}
  \caption{Changes whose human-preferred mapping needs an $N{:}M$ site, per dataset of
  Section~\ref{sec:dataset} and pooled.}
  \label{tab:expressibility}
  \begin{tabular}{lrrr}
    \toprule
    Dataset & $n$ & $N{:}M$ changes & Share \\
    \midrule
    Curated & \AmbiguityCuratedTotal{} & \AmbiguityCuratedWith{} & \AmbiguityCuratedPct{}\% \\
    Full & \AmbiguityFullTotal{} & \AmbiguityFullWith{} & \AmbiguityFullPct{}\% \\
    Stratified\textsuperscript{*} & \AmbiguityStratifiedTotal{} & \AmbiguityStratifiedWith{} &
      \AmbiguityStratifiedPct{}\% \\
    Defects4J\textsuperscript{*} & \AmbiguityDefectsFourJTotal{} & \AmbiguityDefectsFourJWith{} &
      \AmbiguityDefectsFourJPct{}\% \\
    \midrule
    \emph{Curated + Full} & \AmbiguityListScored{} & \AmbiguityListWith{} & \AmbiguityListPct{}\% \\
    \emph{All}\textsuperscript{*} & \AmbiguityScored{} & \AmbiguityAnyFixtures{} &
      \AmbiguityAnyPct{}\% \\
    \bottomrule
  \end{tabular}

  \vspace{2pt}
  \footnotesize\textsuperscript{*}The ambiguity pass over the Stratified and Defects4J changes is
  not finished, so these three figures are lower bounds: a change nobody has reviewed for ambiguity
  is recorded as having none. \emph{Curated + Full} is the rate over the two datasets whose pass is
  complete, and is the figure RA2.3 states.
\end{table}

\begin{resultbox}{RA2.3 (Expressibility)}
  \AmbiguityListPct{}\% of code changes (\AmbiguityListWith{} of \AmbiguityListScored{}) have a
  human-preferred mapping that cannot be expressed at all using the common edit operations
  (Insert, Delete, Update, Move).
\end{resultbox}

This result bears on how tools are scored, because no established tool can produce these
mappings. The node-level accuracy of Section~\ref{sec:omnidiff} counts a group as correct whenever
a tool's mapping is any of the pairings the group allows. The line-level comparison of
Section~\ref{sec:tools} needs one answer per node, so it projects each group through a single
representative pairing, members paired in source order; a tool that chose a different, equally
correct pairing can be charged for the lines on which the two differ.

\textbf{RQ2.4.} What percentage of real-world code changes have a human-preferred mapping whose
cost is higher than that of the cheapest mapping found?

Tree differencing algorithms look for the cheapest mapping, the one with the smallest number of
operations needed to transform one tree into the other. We computed, for each change, the cost of
the human mapping and the cost of a mapping found by a state-of-the-art matcher, under unit costs
for insert, delete and update and a free move. Of the \CostScored{} changes
compared,\textsuperscript{\ref{fn:in-progress}} the two costs are equal \CostTieFixtures{} times.
In \CostHumanHigherFixtures{} (\CostHumanHigherPct{}\%) the human solution is strictly costlier.
The most common reason is that humans treat groups of tokens as one unit, not by their syntax but
by higher-level logic. A common example is the \texttt{return} keyword: if one function was added
and another deleted, matching the two \texttt{return}s lowers the tree edit distance, but a human
finds the match confusing, because the keyword belongs to the logic of the function that contains
it. Existing tools take notice of this: GumTree~\cite{gumtreediff-2014} does not match subtrees
below a threshold size, to avoid such spurious mappings.

\begin{resultbox}{RA2.4 (Cost)}
  On at least \CostHumanHigherPct{}\% of code changes (\CostHumanHigherFixtures{} of
  \CostScored{}), the human-authored mapping costs strictly more, under unit costs for insert,
  delete and update and a free move, than a mapping a tree-differencing algorithm found for the
  same change.
\end{resultbox}

\section{Uniqueness}
\label{sec:uniqueness}

A delta is one of many: any cheapest script will do. This section measures how often the
human-preferred account of a change is not unique either, first in the mapping and then in its
visualisation. Existing tools all provide a single diff. Line-based tools mark lines as inserted
or deleted or match lines one to one, and tree-based tools match nodes one to one. No tool
reports several solutions.

Manually examining real-world changes, we found that some have several equally correct
solutions. Figure~\ref{fig:ambiguous-example} is a motivating example.
\begin{figure}
  \small
  \begin{tabular}{@{}p{0.46\linewidth}@{\hspace{0.04\linewidth}}p{0.46\linewidth}@{}}
    \textsc{Before} & \textsc{After} \\
    \midrule
    \begin{minipage}[t]{\linewidth}
\begin{verbatim}
def is_even(x):
  return x > 0 \
    and x % 2 == 0
\end{verbatim}
    \end{minipage}
    &
    \begin{minipage}[t]{\linewidth}
\begin{verbatim}
def is_even(x):
  return isinstance(x, int) \
    and x > 0 \
    and x % 2 == 0
\end{verbatim}
    \end{minipage} \\
    \centering
    \begin{forest}
      baseline, for tree={font=\ttfamily\scriptsize, inner sep=1.5pt, l sep=5pt, s sep=3pt}
      [return [bool, draw [cmp] [and] [cmp]]]
    \end{forest}
    &
    \centering
    \begin{forest}
      baseline, for tree={font=\ttfamily\scriptsize, inner sep=1.5pt, l sep=5pt, s sep=3pt}
      [return [bool, draw, dashed [bool, draw, dashed [call] [and] [cmp]] [and] [cmp]]]
    \end{forest} \\
  \end{tabular}
  \caption{A change with no unique correct mapping. The before AST holds one
  \texttt{boolean\_operator} node (\texttt{bool}, boxed) and the after AST two, one nested in the
  other (dashed); matching the before node to either costs the same and renders the same.
  \texttt{cmp} is a \texttt{comparison\_operator}.}
  \label{fig:ambiguous-example}
\end{figure}
The AST of the before snippet contains one \texttt{boolean\_operator} node. The after snippet
contains two, one a direct child of the other. There is no objective way to determine which of the
two is added and which one should be matched. The edit distance of both mappings is exactly
equal, and since \texttt{boolean\_operator} is an internal node with no rendering of its own, both
also produce the same visualisation, so human judgement cannot break the tie either.

\textbf{RQ3.1.} What percentage of real-world code changes have multiple optimal node mappings?

We count a change as admitting several optimal mappings when its ground truth records at least one
multi-map group (Section~\ref{sec:ground-truth}). Within a group, any one-to-one pairing that
covers the smaller side is equally correct, and the nodes left over on the other side are deleted
or inserted. Across the dataset the annotators recorded \AmbiguityGroups{} such groups.

Even without such groups, several solutions can exist. Where the annotator's mapping and a
matcher's mapping for one change cost exactly the same under the unit cost model and still
disagree about which nodes pair with which, the cost function cannot separate two different
answers. That happens on \CostTieDifferentFixtures{} of the \CostScored{} changes compared in
Section~\ref{sec:gap} (\CostTieDifferentPct{}\%).

\begin{resultbox}{RA3.1 (Mapping uniqueness)}
  In at least \AmbiguityListPct{}\% of code changes (\AmbiguityListWith{} of
  \AmbiguityListScored{}), humans found that some set of nodes admits several equally correct
  pairings. Independently of that annotation, on \CostTieDifferentPct{}\% of changes the human
  and an algorithmic mapping tie on cost while disagreeing on the mapping, so the cost function
  alone cannot pick between them.
\end{resultbox}

\textbf{RQ3.2.} What percentage of real-world code changes have multiple equally correct
visualisations, and how do those differ?

Settling the mapping does not settle what the reader sees. Of the \PresetPainted{} painted
changes, \PresetBothFixtures{} (\PresetBothPct{}\%) carry a \emph{Minimal} and a \emph{Full}
painting, the annotator's finding that a reader who wants only what changed and a reader who wants
its surroundings marked are both right. \PresetSingleFixtures{} (\PresetSinglePct{}\%) carry one
painting that names no preset: for these a human considers the rendering settled,
whatever the reader's preference. The remaining \PresetMinimalOnlyFixtures{} changes
carry only a \emph{Minimal} painting and \PresetFullOnlyFixtures{} only a \emph{Full} one so far.

../../../plots/preset_table.tex

\begin{figure*}
  \centering
  \includegraphics[width=0.9\linewidth]{figures/preset_difference}
  \caption{The \emph{Full} hand painting against the \emph{Minimal} one, over the changes painted
  both ways. Left: per change, the share of the bytes \emph{Full} paints that \emph{Minimal}
  paints as well, as a cumulative distribution. Right: per language with at least fifteen such
  changes (count in parentheses), the share of the bytes \emph{Full} paints that \emph{Minimal}
  leaves unpainted, by what those bytes hold.}
  \label{fig:presets}
\end{figure*}

We compared the two paintings byte by byte on the \PresetCensusFixtures{} changes that carry
both.\footnote{One change with two paintings is left out: its paintings are two alternatives
that name no preset.} \emph{Minimal} paints \PresetMinimalBytes{} bytes and \emph{Full}
\PresetFullBytes{}, \PresetFullOverMinimal{} times as many. \emph{Full} never leaves unpainted a
byte that \emph{Minimal} paints (\PresetMinimalOnlyBytes{} such bytes), which is one of the
invariants, and the two give \PresetRelabelledBytes{} bytes different labels. The two paintings
are byte-identical on \PresetIdenticalFixtures{} of the changes, as expected: an annotator who
would paint both alike records a single painting instead. Table~\ref{tab:presets} shows what the
\PresetAddedBytes{} bytes that only \emph{Full} paints hold. \PresetAddedWhitespacePct{}\% of
them are whitespace, almost all of it the indentation and spacing of inserted and deleted lines.
\PresetAddedCodePct{}\% are code: the unchanged part of a renamed identifier or a changed number,
which \emph{Full} paints as part of the update, and code that moved, which \emph{Full} marks as
moved in more cases than \emph{Minimal} does. Punctuation alone, the brackets and separators that
a change brings with it, is \PresetAddedPunctuationPct{}\%.

The difference is small for most changes and large for a few. The median change has
\PresetExtraMedian{} bytes that only \emph{Full} paints, the 90th percentile
\PresetExtraPNinety{} and the largest \PresetExtraMax{}. The left panel of
Figure~\ref{fig:presets} shows the distribution relative to the size of the painting: in the
median change, \emph{Minimal} paints \PresetKeptMedianPct{}\% of the bytes \emph{Full} paints.
The right panel shows that the share differs by language, and that what is added differs as
well: mostly whitespace in most languages, mostly code in shell scripts and CSS. We conjecture
that in those two a changed word or value is short and the unchanged remainder that \emph{Full}
paints with it is long by comparison.

\begin{resultbox}{RA3.2 (Visualisation uniqueness)}
  Of the \PresetPainted{} painted changes, \PresetBothFixtures{} (\PresetBothPct{}\%) admit a
  sparse and a generous correct rendering, and \PresetSingleFixtures{} (\PresetSinglePct{}\%)
  have one rendering only. Where both exist, the generous one paints \PresetFullOverMinimal{}
  times the bytes, and \PresetAddedWhitespacePct{}\% of what it adds is whitespace.
\end{resultbox}

For Section~\ref{sec:background} this is the last piece of evidence: even the visualisation of a
change is a family, with the reader's preference as its parameter, and the family has a single
member for \PresetSinglePct{}\% of the painted changes. A delta has no such parameter.

\section{Exact Deltas within an Interactive Budget}
\label{sec:speed}

This section asks what the exact delta costs to compute. The data is for code, the content whose
exact delta is expensive: a line script or a pixel difference is computed in time that is close
to linear for typical inputs. The asymptotic complexity of APTED and Zhang--Shasha is $O(n^3)$
for general trees, where $n$ is the number of nodes. APTED performs considerably better on real
trees because it is less sensitive to the tree shape. Given this direction in the literature, we
focus on source-code trees and ask the following question:

\textbf{RQ4.} What percentage of real-world source-code changes can a single, whole-tree
tree-edit-distance computation complete within one second?

One second was chosen as a budget because it is where interaction stops feeling immediate.
Nielsen's response-time limits~\cite{nielsen-usability} put one second at the boundary at which a
user's flow of thought stays uninterrupted but the delay becomes noticeable; below about 100\,ms
a response reads as instantaneous. Practitioners rate the reaction time of review tooling as very
important~\cite{kudrjavets-2023}, and developer satisfaction with code review falls as changes
wait longer~\cite{meta-review-time-2022}. Developers have had decades to get used to instant
diffs from line-based tools, so any tool that expects to be adopted has to meet a similar bar.

We sampled \RqOnePairsAttempted{} real (before, after) file pairs from single-parent commits in
the corpus, stratified per language and per size bucket, with the buckets defined in
Section~\ref{sec:diff-shape}: up to \RqOnePerCell{} pairs in each of the \RqOneCells{} cells of
\RqOneLanguages{} languages and \RqOneBuckets{} buckets. \RqOneFullCells{} cells reached that
target; the rest hold every pair the corpus had. On each pair, we ran our optimised APTED
implementation, with a one-second budget enforced by operating-system limits.

\begin{figure*}
  \centering
  \includegraphics[width=\linewidth]{figures/time_budget}
  \caption{The one-second budget, read the same way twice. Left: the share of sampled file pairs
  whose whole-tree APTED computation completes within the time on the axis, per artifact
  category; the height of each curve at the dashed line is RA4, and the pairs the budget killed
  are the gap to 100\%. Right: wall-clock per diff for every tool configuration of
  Section~\ref{sec:tools}, every repeated run over every fixture, with \textsc{OmniDiff} charged
  for its parse and drawn dash-dotted; the dotted curve is the tree-sitter parse alone, the floor
  every AST-aware tool pays first. The dashed line is the same one second in both panels.}
  \label{fig:time-budget}
\end{figure*}

The left panel of Figure~\ref{fig:time-budget} shows, for each artifact category, the share of
pairs completed as a function of the time allowed. Size decides the outcome: within one second
the computation finishes for \RqOneCodeThirtyToHundredPct{}\% of code file pairs at 31 to 100
lines and for \RqOneCodeHundredToThreeHundredPct{}\% at 101 to 300. The shape of the input
matters as well: scripting languages complete more often than code, and data files less often.

\begin{resultbox}{RA4 (Exactness)}
  Within one second, a whole-tree tree-edit-distance computation finds the delta for only
  \RqOneCodePct{}\% of code file pairs.
\end{resultbox}

\section{State of the Art}
\label{sec:tools}

This section measures what established tools do with real changes: what they show for each type
of content, and, for code, how accurately and how fast.

%% RQ5.1-RQ5.3 are stated here, next to the measurements that answer them. Section 1 names them
%% by group; keep the wording in step.
\textbf{RQ5.1.} What do established tools show for each type of content?

For text, \texttt{diff} and \texttt{git diff} show a line script, and they do so for every text
file, whatever it holds. For the \ChangeCensusBinaryPct{}\% of changed files that git judges
binary (RA1.1), git shows that the two versions differ and nothing else, unless a user has
configured a converter for the format. The tools that go further each cover a part of the
content. The AST-aware tools read code in the languages they have a grammar for: of the
\UnixDiffFixtures{} code changes of the comparison below, diffsitter's compiled-in grammar set
processes \DiffsitterFixtures{}, GumTree's generator set \GumTreeFixtures{}, difftastic's
\DifftasticFixtures{}, and srcDiff, which parses five languages, \SrcDiffFixtures{}. diffoscope
unpacks containers and compares what is inside them as text or as bytes, and pixelmatch counts
the pixels of two pictures that differ visibly. None of the tools we know of states for a
changed picture, font or archive how much of the change a viewer would see, which is what the
verdicts of Section~\ref{sec:ground-truth} record, so we could score no established tool against
them.

\begin{resultbox}{RA5.1 (Coverage)}
  Git shows a line script for the \ChangeCensusTextPct{}\% of changed files it reads as text and
  only a notice for the \ChangeCensusBinaryPct{}\% it judges binary. Of four AST-aware tools, the
  broadest reads \DifftasticFixtures{} of \UnixDiffFixtures{} sampled code changes and the
  narrowest \SrcDiffFixtures{}.
\end{resultbox}

\textbf{RQ5.2.} How accurately do current state-of-the-art code diffing tools find the
human-preferred visualisation?

\textbf{RQ5.3.} How fast are they?

\textbf{Comparison tools.} We measure eight established tools, all introduced in
Section~\ref{sec:background-tools}, on \UnixDiffFixtures{} of the code changes of
Section~\ref{sec:dataset}.\textsuperscript{\ref{fn:in-progress}} \texttt{git diff} is measured
under each of its four algorithms, Myers, minimal, patience and histogram, giving eleven
configurations in all. Five configurations report at line granularity only: Unix \texttt{diff}
and the four \texttt{git} algorithms. Six report sub-line detail: BDiff~\cite{bdiff}, Neovim's
diff mode, and the four AST-aware tools, GumTree (\GumTreeVersion{})~\cite{gumtreediff-2014},
difftastic~\cite{difftastic}, diffsitter~\cite{diffsitter} and srcDiff~\cite{decker-2020}. Our
own tool's figures are reported in Section~\ref{sec:omnidiff}, on the same basis and against the
same ground truth.

\textbf{Line-level agreement.} None of the eleven configurations reports a node mapping or a
painting in a form that can be compared directly, so we compare them on a common ground: which
source lines each marks as touched, against the lines the human mapping touches. Because the
tools cover different subsets of the changes (RA5.1), percentages are used.

We report agreement per change rather than pooled over lines in
Table~\ref{tab:agreement-buckets}, and Figure~\ref{fig:agreement-by-dataset} splits the same
reading by the dataset each change was sampled from. \emph{Perfect} means zero mismatched lines.
$\ge$99\% is ``a developer has to look twice''; 95--99\% is ``a developer has to check''. Below
95\% the tool is no longer usable.

\begin{resultbox}{RA5.2 (Tool accuracy)}
srcDiff renders \SrcDiffPerfectPct{}\% of the code changes it parses perfectly, but it parses
five languages; no tool that parses most of the changes renders much more than three fifths of
them perfectly, the best being difftastic at \DifftasticPerfectPct{}\%.
\end{resultbox}

\input{figures/benchmark_other_buckets_combined.tex}

\begin{figure*}
  \centering
  \includegraphics[width=\linewidth]{figures/benchmark_other_buckets_by_dataset}
  \caption{The line-level agreement of Table~\ref{tab:agreement-buckets}, split by the dataset
  of Section~\ref{sec:dataset} each fixture was drawn from, with the pooled reading on the right.
  Each bar is one configuration over the fixtures it supports ($n$ beside the bar); the segments
  are the four agreement buckets, and the Perfect segment carries its share.}
  \label{fig:agreement-by-dataset}
\end{figure*}

The split by dataset in Figure~\ref{fig:agreement-by-dataset} orders the datasets the same way for
every configuration: each scores its own best on either the Stratified sample or the Defects4J
changes, and all five line-granularity configurations do best on Defects4J. Bug fixes are short,
localised edits, which is the shape a line-based diff recovers most reliably, so the ordering is a
property of the datasets rather than of any tool.

\textbf{Cost.} Table~\ref{tab:speed-vs-others} shows wall-clock percentiles, sorted fastest to
slowest by median. Each tool is measured over the changes it supports rather than over one common
subset, so the populations behind these rows differ: the text-based tools draw on every change,
while an AST tool draws only on the languages its grammars cover. Two tools are reported twice:
GumTree once invoked as a fresh process per change, the realistic cost of a single command-line
invocation, and once run inside one persistent JVM across all changes, isolating JVM startup from
GumTree's own matching cost; and BDiff likewise per process and inside one persistent Python
interpreter, for the same reason.

\begin{table}
  \caption{Wall-clock time percentiles for the eleven tool configurations, milliseconds, sorted
  fastest to slowest by median, with \textsc{OmniDiff} on the last row. \emph{Max} is the slowest
  single one of the \TimingRepeats{} runs over any fixture. Every configuration is timed as a whole
  process except \textsc{OmniDiff}, which is timed in-process and so pays no process startup; see
  Section~\ref{sec:threats}.}
  \label{tab:speed-vs-others}
  \small
  \begin{tabular}{lrrrr}
    \toprule
    Tool & p50 & p90 & p99 & Max \\
    \midrule
    \texttt{git} (histogram) & \SpeedGitHistogramPFifty{} & \SpeedGitHistogramPNinety{} & \SpeedGitHistogramPNinetyNine{} & \SpeedGitHistogramMax{} \\
    \texttt{git} (minimal) & \SpeedGitMinimalPFifty{} & \SpeedGitMinimalPNinety{} & \SpeedGitMinimalPNinetyNine{} & \SpeedGitMinimalMax{} \\
    Unix \texttt{diff} & \SpeedUnixDiffPFifty{} & \SpeedUnixDiffPNinety{} & \SpeedUnixDiffPNinetyNine{} & \SpeedUnixDiffMax{} \\
    \texttt{git} (patience) & \SpeedGitPatiencePFifty{} & \SpeedGitPatiencePNinety{} & \SpeedGitPatiencePNinetyNine{} & \SpeedGitPatienceMax{} \\
    \texttt{git} (Myers) & \SpeedGitMyersPFifty{} & \SpeedGitMyersPNinety{} & \SpeedGitMyersPNinetyNine{} & \SpeedGitMyersMax{} \\
    BDiff (warm interpreter) & \SpeedBDiffWarmPFifty{} & \SpeedBDiffWarmPNinety{} & \SpeedBDiffWarmPNinetyNine{} & \SpeedBDiffWarmMax{} \\
    diffsitter & \SpeedDiffsitterPFifty{} & \SpeedDiffsitterPNinety{} & \SpeedDiffsitterPNinetyNine{} & \SpeedDiffsitterMax{} \\
    GumTree (warm JVM) & \SpeedGumTreeWarmPFifty{} & \SpeedGumTreeWarmPNinety{} & \SpeedGumTreeWarmPNinetyNine{} & \SpeedGumTreeWarmMax{} \\
    srcDiff & \SpeedSrcDiffPFifty{} & \SpeedSrcDiffPNinety{} & \SpeedSrcDiffPNinetyNine{} & \SpeedSrcDiffMax{} \\
    \texttt{nvim -d} & \SpeedNvimDiffPFifty{} & \SpeedNvimDiffPNinety{} & \SpeedNvimDiffPNinetyNine{} & \SpeedNvimDiffMax{} \\
    difftastic & \SpeedDifftasticPFifty{} & \SpeedDifftasticPNinety{} & \SpeedDifftasticPNinetyNine{} & \SpeedDifftasticMax{} \\
    BDiff (cold, per-invocation) & \SpeedBDiffColdPFifty{} & \SpeedBDiffColdPNinety{} & \SpeedBDiffColdPNinetyNine{} & \SpeedBDiffColdMax{} \\
    GumTree (cold, per-invocation) & \SpeedGumTreeColdPFifty{} & \SpeedGumTreeColdPNinety{} & \SpeedGumTreeColdPNinetyNine{} & \SpeedGumTreeColdMax{} \\
    \midrule
    \textsc{OmniDiff} & \SpeedOmniDiffPFifty{} & \SpeedOmniDiffPNinety{} & \SpeedOmniDiffPNinetyNine{} & \SpeedOmniDiffMax{} \\
    \bottomrule
  \end{tabular}
\end{table}

\begin{resultbox}{RA5.3 (Tool cost)}
The eleven configurations span more than two orders of magnitude at the median, from Unix
\texttt{diff}'s \SpeedUnixDiffPFifty{}\,ms to GumTree's \SpeedGumTreeColdPFifty{}\,ms per
invocation. At the 99th percentile the same range runs from \SpeedUnixDiffPNinetyNine{}\,ms to
\SpeedGumTreeColdPNinetyNine{}\,ms, three orders of magnitude apart.
\end{resultbox}

We report both a cold and a warm figure for GumTree and BDiff because they answer different
questions: the cold number is what a developer running a diff from the command line waits for,
and the warm number is the matching cost once process startup is amortised away. Every timing
number in this paper is the result of repeated runs rather than a single measurement. Each tool
is run \TimingRepeats{} times on every change, and the percentiles are taken over every
individual run pooled across changes rather than over per-change means: the repeat count is
uniform, so no change is over-weighted.

\section{\textsc{OmniDiff}}
\label{sec:omnidiff}

This section is the paper's tool contribution. \textsc{OmniDiff} diffs a file by what it holds.
It recognises the content of a pair from its bytes (Section~\ref{sec:omnidiff-recognition}),
diffs code with a grammar by mapping its AST and projecting the mapping to a painting, and any
other text line by line (Section~\ref{sec:omnidiff-code}), and diffs binary content by one of
two content engines (Section~\ref{sec:omnidiff-content}). Table~\ref{tab:findings-to-design}
shows what each empirical result decided in its design, and Section~\ref{sec:omnidiff-evaluation}
evaluates the tool against the same ground truth and on the same basis as the tools of
Section~\ref{sec:tools}.

\begin{table}
  \caption{What each empirical result decided in \textsc{OmniDiff}'s design.}
  \label{tab:findings-to-design}
  \small
  \begin{tabular}{@{}>{\raggedright\arraybackslash}p{0.44\linewidth}>{\raggedright\arraybackslash}p{0.52\linewidth}@{}}
    \toprule
    Finding & Consequence \\
    \midrule
    \ChangeCensusBinaryPct{}\% of changed files are binary to git, and pictures change in
      \ContentCensusPicturesRepos{} repositories (RA1.1)
      & Content is recognised by its bytes and diffed by what it holds \\ \addlinespace
    The median edit rewrites \EditsChurnPFiftyPct{}\% of its file (RA1.2)
      & Phase 1 settles the unchanged majority by hashing before any edit-distance computation
        runs \\ \addlinespace
    A tree script with the right mapping breaks a painting invariant on
      \InvariantTreeScriptFixturesPct{}\% of changes (RA2.1)
      & The renderer is a projection with rules of its own, not a colouring of the mapping \\ \addlinespace
    \InvisibleNodeShare{}\% of nodes never reach the screen, and \ContentVerdictNothingPct{}\%
      of content changes show nothing deliberate (RA2.2)
      & Accuracy is reported on visible nodes and painted bytes; every content change is named
        by how much of it a viewer sees \\ \addlinespace
    \AmbiguityListPct{}\% of code changes have no expressible optimum and
      \CostTieDifferentPct{}\% tie at different mappings (RA2.3, RA3.1)
      & The phase order is the tie-break: the earliest phase to claim a node wins, and the
        closing step marks what no phase claimed \\ \addlinespace
    On \CostHumanHigherPct{}\% of changes the human mapping costs more than the cheapest one
      (RA2.4)
      & Phases 2 and 3 encode reader preferences that cost more than the optimum: comments and
        modifiers follow their node, names beat cost, diagnostics are held back \\ \addlinespace
    \PresetBothPct{}\% of painted changes have a sparse and a generous rendering (RA3.2)
      & Two render presets, \emph{Minimal} and \emph{Full}, over one mapping \\ \addlinespace
    Whole-tree APTED finishes within one second for only \RqOneCodePct{}\% of code pairs (RA4)
      & APTED runs only inside one already-matched pair (Phase 3), never over the whole tree \\ \addlinespace
    No established tool that parses most code changes renders more than
      \DifftasticPerfectPct{}\% of them perfectly (RA5.2)
      & The bar the evaluation has to clear, on the same basis \\ \addlinespace
    Files reach \AstMax{} nodes (Section~\ref{sec:files})
      & The Robust target below \\
    \bottomrule
  \end{tabular}
\end{table}

\textsc{OmniDiff} has two targets:

\begin{itemize}
  \item \textbf{Robust:} \textsc{OmniDiff} must handle every file up to the measured maximum,
    \AstMax{} AST nodes, with headroom above that number.
  \item \textbf{Fast:} \textsc{OmniDiff} must compute a diff in under \SpeedTargetMs{}\,ms for
    \SpeedTargetPct{}\% of real-world changes, RQ4's one second applied to the population of
    real changes rather than to a sample of hard cases.
\end{itemize}

\subsection{Recognising content}
\label{sec:omnidiff-recognition}

A content kind is recognised from the file's bytes, never from its name: under \texttt{git diff}
both sides are temporary files, and an X cursor is an extensionless file named after the cursor.
A pair is diffed as one kind or not at all, and a pair that does not decode is reported as any
other binary pair rather than as an error, so that, run as git's external diff, one bad file does
not abandon every file after it. Text needs no such recognition, but two kinds of it are widened:
XML formats such as SVG are parsed with the XML grammar, and text in UTF-16 or UTF-32 with a byte
order mark is decoded and diffed as text rather than as binary. Text in one of
\NumSupportedLanguages{} languages with a tree-sitter grammar~\cite{tree-sitter} is diffed as
code; any other text is diffed line by line, with the changed part of a paired line marked.

\subsection{Text and code}
\label{sec:omnidiff-code}

For code, \textsc{OmniDiff} first computes a mapping and then projects it to a painting, the two
steps of Section~\ref{sec:definitions}. It matches two ASTs in \NumPhasesWord{} phases. Each
phase either resolves part of the mapping directly or narrows the residual left for the next
phase. Later phases only ever see nodes earlier phases left unmatched. A sixth step closes the
mapping, recording the deletion or insertion implied by whatever every phase above left
undecided.

\textbf{Phase 1, hash-based descent.} \textsc{OmniDiff} matches subtrees by two hash variants, in
order: byte-identical subtrees, then same-shape subtrees with any leaf value. This phase is
motivated directly by RA1.2: since most commits touch a small fraction of a file's lines, most
of a typical file's AST is byte-identical between before and after. Resolving that identical
majority with cheap hashing, before any tree-edit-distance computation runs at all, avoids paying
$O(n^3)$ on content that never changed.

\textbf{Phase 2, contextual exact matching.} Comments and modifiers (attributes, decorators) that
immediately precede an already-matched node, and diagnostic statements (logging calls, assertions,
panics) that are byte-identical on both sides, get matched directly. Neither is reachable by
phase 1's structural hashing: comments and modifiers are not part of the hashed AST shape, and
diagnostic statements are held back deliberately, so that matching one does not fragment a larger
match around it. These are the preferences RA2.4 measures: each of them can cost more than the
cheapest mapping.

\textbf{Phase 3, syntax-aware subtree matching.} This phase generalises three matching strategies:
matching by fully-resolved name (handling overloads and duplicate names across scopes), a greedy
cost-estimate matcher for containers with no name structure, such as an \texttt{if} block or a
loop body, and an overlap matcher pairing import statements that share a base path. Before any of
these run, large flat subtrees, such as a macro's token list or a big flat argument list, are
pre-matched directly with Myers' algorithm~\cite{myers-1986}, since a flat sequence has no tree
structure worth exploiting. Each accepted pair is then handed to APTED~\cite{apted-2015}, scoped
to that pair alone, to compute its exact internal edit script. This is RA4 applied: the exact
computation is affordable inside one matched pair, and only there.

\textbf{Phase 4, residual alignment.} Whatever remains unmatched after phases 1 through 3 is
resolved by a Myers-LCS alignment~\cite{myers-1986} over the residual forest, bracketed by two
passes of strict bottom-up propagation: a node whose children all match the children of exactly
one node on the other side is matched to it, which shrinks the residual, and is run again
afterward so that pockets the alignment just resolved can still bubble up to their now-complete
ancestors.

\textbf{Phase 5, move-detection recovery.} The final matching pass. Ordered tree edit distance
cannot express a subtree that relocated across a matched boundary as anything but a delete plus
an insert. This phase, adapted directly from GumTree's recovery-mappings
phase~\cite{gumtreediff-2014}, pairs up whole subtrees left fully deleted on one side and fully
inserted on the other, when they are byte-identical and large enough to rule out coincidence.

\textbf{Projection.} The renderer turns the mapping into a painting under one of two presets,
\emph{Minimal} and \emph{Full}, the two the ground truth is painted in. It is not the literal
projection: following RA2.1, it decides separately which whitespace, separators and brackets go
with a change, whether a relocated node is shown as moved or left unmarked as a pure reindent,
and how much of a changed word to mark. Its rules follow the painting invariants of
Section~\ref{sec:invariants}, and where they can the two share code: the list of tokens that are
painted whole is read by the renderer and by the invariant alike.

\textbf{Code inside strings.} After the five phases, every matched pair of host strings that
changed is examined for code of another language. Where the string's context names a language
(a sibling \texttt{"path"} naming a file, a YAML \texttt{run:} key, a \texttt{css} template tag)
and both sides parse without an error, the decoded strings are diffed by the whole pipeline as
files of their own, and the result is painted back onto the host file through a map from each
decoded byte to its source. A changed string whose language is not detected is still diffed as
text, line by line, so two edits far apart in one long string are painted apart rather than as
one span from the first change to the last.

\subsection{Other content}
\label{sec:omnidiff-content}

\textbf{Two engines.} A picture is pixels. Which pixels changed is decided by pixelmatch's
perceptual test~\cite{pixelmatch}, so that re-encoding noise is not reported as change; changed
pixels close to each other are grouped into one region, and the frames of an animation are aligned
the way a text diff aligns lines. Everything else is a container of named members: an archive's
files, a font's glyphs, a PDF's pages, a catalog's messages, a cursor theme's sizes. Members match
by key and are hashed by what they show, so only a changed member is decoded, and it is diffed by
what it holds in turn, a picture in a zip as a picture. A glyph is keyed by the character it draws
and hashed by its outline and advance width, so a font re-hinted or re-wrapped as WOFF2 has no
changed glyph. A PDF page is keyed by its number and drawn as a picture, a message by what it
translates, and a Java class file in a jar is listed as its fields and methods and what each
method's bytecode refers to, which diffs as text.

\textbf{Verdicts.} Each change, and each changed member of a container, is named by one of the
six levels of Section~\ref{sec:ground-truth}. This is RA2.2 applied to content: the byte delta of
such a pair is never empty, and a diff that reports a lossless re-encoding the way it reports a
redrawn icon hides the one among the other.

\subsection{Evaluation}
\label{sec:omnidiff-evaluation}

\textbf{Coverage of content.} Of the \ContentCensusChanges{} changed files of RA1.1,
\textsc{OmniDiff} diffs \ContentCensusShownPct{}\% as text or by their content, against the
\ChangeCensusTextPct{}\% git diffs as text. Counting at most \ContentCensusCapPerRepo{} changes
per repository, so that a few very large repositories do not decide it, the share is
\ContentCensusShownCappedPct{}\%, and \ContentCensusReposEntirely{} of the \ContentCensusRepos{}
repositories are diffed entirely.

\begin{table}
  \caption{\textsc{OmniDiff} against the code ground truth, per change. \emph{Exact}: no visible
  node mapped differently from the human mapping, or no byte painted differently from the hand
  painting. \emph{Within 1\%}: at most 1\% of the visible nodes, or of the bytes of the two
  files. \emph{Human mapping}: exact when the renderer is given the human mapping.}
  \label{tab:omnidiff-accuracy}
  \small
  \setlength{\tabcolsep}{4pt}
  \begin{tabular}{@{}lrrrr@{}}
    \toprule
    Compared & $n$ & Exact & Within 1\% & Human mapping \\
    \midrule
    Node mapping & \MappingFixtures{} & \MappingZeroVisiblePct{}\% & \MappingWithinOnePct{}\% & -- \\
    Painting, \emph{Minimal} & \PaintingMinimalFixtures{} & \PaintingMinimalZeroPct{}\% & \PaintingMinimalWithinOnePct{}\% & \PaintingMinimalIdealZeroPct{}\% \\
    Painting, \emph{Full} & \PaintingFullFixtures{} & \PaintingFullZeroPct{}\% & \PaintingFullWithinOnePct{}\% & \PaintingFullIdealZeroPct{}\% \\
    Painting, one rendering & \PaintingSingleFixtures{} & \PaintingSingleZeroPct{}\% & \PaintingSingleWithinOnePct{}\% & \PaintingSingleIdealZeroPct{}\% \\
    \bottomrule
  \end{tabular}
\end{table}

\textbf{Code: mapping.} Against the human mapping, \textsc{OmniDiff} maps \NodeAccuracyPct{}\% of
AST nodes correctly (\NodeMismatches{} mismatches in \NodesTotal{}), and
\VisibleNodeAccuracyPct{}\% of the \VisibleNodeShare{}\% of nodes that carry text and can reach
the screen (\VisibleNodeMismatches{} in \VisibleNodesTotal{}). Per change,
Table~\ref{tab:omnidiff-accuracy} shows that it maps every visible node as the human did on
\MappingZeroVisiblePct{}\% of the \MappingFixtures{} changes (\MappingZeroVisible{}).

\textbf{Code: painting.} The painting is compared byte by byte with the hand painting, in the
three cases of Section~\ref{sec:ground-truth}; where a change has one rendering, both presets
have to produce it, and a byte counts as wrong when either does not.
Table~\ref{tab:omnidiff-accuracy} shows the result. \textsc{OmniDiff} paints every byte as the
human did on \PaintingMinimalZeroPct{}\% of the changes under \emph{Minimal}, on
\PaintingFullZeroPct{}\% under \emph{Full}, and on \PaintingSingleZeroPct{}\% of the changes with
one rendering. The last column separates the two steps: given the human mapping instead of its
own, the renderer is exact on \PaintingMinimalIdealZeroPct{}\%, \PaintingFullIdealZeroPct{}\% and
\PaintingSingleIdealZeroPct{}\%. Most of what is left is therefore the projection's and not the
matcher's, and it is largest where RA3.2 puts the reader's preference: in what \emph{Full} adds.

\textbf{Code: invariants.} We also hold \textsc{OmniDiff}'s output against the
\InvariantLabelCount{} invariants that can be checked on a painting's labels alone, on the
\InvariantFixtures{} changes of RA2.1. Its rendering of its own mapping breaks one of them on
\InvariantLabelToolRenderedFixturesPct{}\% of the changes, and its rendering of the human mapping
on \InvariantLabelHumanRenderedFixturesPct{}\%. By the same rules the tree script, which is that
same human mapping under the literal projection, breaks one on
\InvariantLabelTreeScriptFixturesPct{}\%, the line script on
\InvariantLabelLineScriptFixturesPct{}\% and the hand paintings on
\InvariantLabelHumanFixturesPct{}\%. The projection thus closes most of the distance between a
correct mapping and a visualisation, and the matcher's errors account for the difference between
the first two figures.

\textbf{Code: comparison with established tools.} On the line-level basis of
Section~\ref{sec:tools} and its \OmniDiffFixtures{} changes,\textsuperscript{\ref{fn:in-progress}}
\textsc{OmniDiff} renders \OmniDiffPerfectPct{}\% of the changes perfectly, against
\DifftasticPerfectPct{}\% for the best established tool that parses most of them, and it
completes every one. It marks \OmniDiffLineMismatches{} lines differently from the human mapping,
a rate of \OmniDiffLineRate{}\%, the lowest of any tool measured. On the \SrcDiffFixtures{}
changes srcDiff parses, it renders \SrcDiffSubsetOmniDiffPerfectPct{}\% perfectly against
srcDiff's \SrcDiffPerfectPct{}\%. Table~\ref{tab:omnidiff-by-dataset} splits the result by
dataset. A line is a coarser unit than a painted byte, which is why this figure is higher than
those of Table~\ref{tab:omnidiff-accuracy}.

../../../plots/benchmark_omnidiff_by_dataset.tex

\textbf{Code: robustness.} Run over the \RobustnessFixtures{} fixtures with a
\RobustnessTimeoutSeconds{}-second budget per fixture and no cap on tree size,
\textsc{OmniDiff} completes every one. The largest fixture holds \RobustnessMaxNodes{}
AST nodes per side, the slowest takes \RobustnessSlowestMs{}\,ms, and peak heap use is
\RobustnessPeakMemoryMb{}\,MB. The same measurement over the whole Full corpus reaches the files
the fixtures do not. Every modified code file in the corpus's recent history
(Section~\ref{sec:content}), \RobustnessFullAttempted{} pairs across
\RobustnessFullLanguages{} languages, was diffed with no node cap under a
\RobustnessFullTimeoutSeconds{}-second budget and a \RobustnessFullMemoryCapGb{}\,GB memory cap
per process. \textsc{OmniDiff} completes \RobustnessFullCompletedPct{}\% of them with
\RobustnessFullPanicked{} panics, at a median of \RobustnessFullPFiftyMs{}\,ms and a 99th
percentile of \RobustnessFullPNinetyNineMs{}\,ms; the largest completed pair holds
\RobustnessFullMaxSideNodes{} AST nodes on each side, against the \AstMax{} of the largest file
in the corpus that the Robust target names. The remainder is \RobustnessFullTimedOut{} pairs past
the budget and \RobustnessFullKilled{} past the memory cap. The killed pairs are commits of
\RobustnessFullKilledFiles{} distinct files, all but one of them generated parser tables,
minified bundles or embedded binary data; the exception is one commit of a 40,000-line
single-header C++ library. Given \RobustnessFullBigMemoryCapGb{}\,GB and a five-minute budget, one
commit of each memory-killed file was tried again: \RobustnessFullKilledFilesCompletedUnderBigCap{}
completed, one needed more memory still, and the rest ran past even that five-minute budget.

\textbf{Code: speed.} Over the \OmniDiffFixtures{} changes of the tool comparison,
\textsc{OmniDiff}'s median wall-clock time is \SpeedOmniDiffPFifty{}\,ms, its 90th percentile
\SpeedOmniDiffPNinety{}\,ms and its 99th percentile \SpeedOmniDiffPNinetyNine{}\,ms, parse
included. The right panel of Figure~\ref{fig:time-budget} places that curve among the eleven
configurations of Table~\ref{tab:speed-vs-others}: at the median the line-based tools sit to its
left and the AST-aware tools to its right, and besides diffsitter, over the \DiffsitterFixtures{}
changes its grammars cover, it is the only AST-aware curve that reaches 100\% before the budget
line.

\begin{table}
  \caption{Speed and accuracy per content family. Speed is milliseconds per sampled pair,
  decoding included; accuracy is against the human verdicts on the judged changes.}
  \label{tab:content}
  \small
  \setlength{\tabcolsep}{3.5pt}
  \begin{tabular}{@{}lrrrrrrr@{}}
    \toprule
    & \multicolumn{4}{c}{Speed} & \multicolumn{3}{c}{Accuracy} \\
    \cmidrule(lr){2-5}\cmidrule(l){6-8}
    Content & $n$ & p50 & p90 & p99 & $n$ & Perfect & Near \\
    \midrule
    Pictures & \ContentSpeedPicturesPairs{} & \ContentSpeedPicturesPFifty{} & \ContentSpeedPicturesPNinety{} & \ContentSpeedPicturesPNinetyNine{} & \ContentAccPicturesJudged{} & \ContentAccPicturesPerfectPct{}\% & \ContentAccPicturesNearPerfectPct{}\% \\
    Archives & \ContentSpeedArchivesPairs{} & \ContentSpeedArchivesPFifty{} & \ContentSpeedArchivesPNinety{} & \ContentSpeedArchivesPNinetyNine{} & \ContentAccArchivesJudged{} & \ContentAccArchivesPerfectPct{}\% & \ContentAccArchivesNearPerfectPct{}\% \\
    PDFs & \ContentSpeedPdfsPairs{} & \ContentSpeedPdfsPFifty{} & \ContentSpeedPdfsPNinety{} & \ContentSpeedPdfsPNinetyNine{} & \ContentAccPdfsJudged{} & \ContentAccPdfsPerfectPct{}\% & \ContentAccPdfsNearPerfectPct{}\% \\
    Fonts & \ContentSpeedFontsPairs{} & \ContentSpeedFontsPFifty{} & \ContentSpeedFontsPNinety{} & \ContentSpeedFontsPNinetyNine{} & \ContentAccFontsJudged{} & \ContentAccFontsPerfectPct{}\% & \ContentAccFontsNearPerfectPct{}\% \\
    Catalogs & \ContentSpeedCatalogsPairs{} & \ContentSpeedCatalogsPFifty{} & \ContentSpeedCatalogsPNinety{} & \ContentSpeedCatalogsPNinetyNine{} & \ContentAccCatalogsJudged{} & \ContentAccCatalogsPerfectPct{}\% & \ContentAccCatalogsNearPerfectPct{}\% \\
    Cursors & \ContentSpeedCursorsPairs{} & \ContentSpeedCursorsPFifty{} & \ContentSpeedCursorsPNinety{} & \ContentSpeedCursorsPNinetyNine{} & \ContentAccCursorsJudged{} & \ContentAccCursorsPerfectPct{}\% & \ContentAccCursorsNearPerfectPct{}\% \\
    \midrule
    \emph{All} & \ContentSpeedPairs{} & \ContentSpeedPFifty{} & & & \ContentAccAllJudged{} & \ContentAccAllPerfectPct{}\% & \ContentAccAllNearPerfectPct{}\% \\
    \bottomrule
  \end{tabular}
\end{table}

\textbf{Other content: speed.} Table~\ref{tab:content} shows speed and accuracy per content
family. Speed is timed over \ContentSpeedPairs{} pairs sampled from the Full list, on the machine
of Section~\ref{sec:project}, one diff per pair from the bytes in memory, decoding included:
unpacking an archive, drawing a font's glyphs or a PDF's pages. Pictures, catalogs and cursors meet
the Fast target, with \ContentSpeedPicturesUnderSecondPct{}\% of pictures under one second.
Archives, fonts and PDFs do not: \ContentSpeedArchivesUnderSecondPct{}\% of archives,
\ContentSpeedFontsUnderSecondPct{}\% of fonts and \ContentSpeedPdfsUnderSecondPct{}\% of PDFs
finish within one second, because each changed member is decoded in turn and a font or a PDF
draws every changed glyph or page. Over all families, \ContentSpeedUnderSecondPct{}\% of pairs
finish within one second, against the \SpeedTargetPct{}\% the target asks.

\textbf{Other content: accuracy.} \textsc{OmniDiff}'s verdicts equal every human verdict on
\ContentAccAllPerfectPct{}\% of the \ContentAccAllJudged{} judged changes
(\ContentAccAllPerfect{}), and are all within one level on \ContentAccAllNearPerfectPct{}\%.
Pictures, the largest set, are the weakest, at \ContentAccPicturesPerfectPct{}\% perfect. We
conjecture that most of that gap lies in the picture engine's thresholds between imperceptible,
artifacts, edited and redrawn, which are set by hand and not yet fitted to the judged pictures. The
PDF figure rests on \ContentAccPdfsJudged{} changes.

\section{Threats to Validity}
\label{sec:threats}

\textbf{Construct validity.} The invariants were written by the people who made the paintings,
and the paintings were then repaired to obey them (Section~\ref{sec:invariants}). The low rate
at which the hand paintings break them is therefore partly by construction, and RA2.1 measures
the distance of a script from this convention of visualisation, not from every possible one. Two
choices in that measurement favour or disfavour the scripts and are stated with it: the tree
script is given the human mapping, which no algorithm produces, and both scripts are filed under
both presets, so that the line script is charged for indentation under \emph{Minimal}. The share
of changes on which the line script marks surviving code as removed does not depend on the second
choice.

Every rate this paper reports about ambiguity is a floor rather than an estimate. A multi-map
group exists where an annotator noticed that two pairings were equally correct and recorded it; a
change whose alternative mapping nobody saw is filed as unique. The same holds for the paintings:
a change with a single painting is one for which the annotator saw one rendering. RA3.1's second
reading, cost ties at differing mappings, needs no annotation to detect and exists for that
reason. We score RA2.3 and RA3.1 over the \AmbiguityListScored{} changes whose ambiguity pass is
complete rather than over all \NumFixtures{}. That pass covers the Curated and Full lists in
full, so for those two lists the bound is one on annotator attention rather than one on coverage.
The Stratified sample (\AmbiguityStratifiedWith{} of \AmbiguityStratifiedTotal{} changes) runs
far under the rate the two finished lists show, which is an annotation gap rather than a finding
about those changes; Table~\ref{tab:expressibility} prints its row and the \emph{All} row under
that caveat.

Because no whole-file exact computation is affordable at these sizes (RA4), RA2.4 compares the
cost of the human mapping to that of a heuristic computation rather than to a proven optimum. Its
rate is therefore a lower bound.

Section~\ref{sec:tools} scores tools at the line level because that is the only granularity all
eleven configurations share. A line-level projection of a node mapping discards detail that the
AST-aware tools compute and the line-based ones never had, which if anything favours the
line-based tools. It also means the paintings never enter that comparison, so RA3.2's error term
sits in those tables unremoved.

The content ground truth is a verdict, not a painting. A matching verdict says that
\textsc{OmniDiff} judged how much a change changed as the annotator did, not that it found where:
the regions and members it reports are not scored, except that a member the annotator judged has
to be among them. The six levels are also coarser than a node mapping, so a perfect content
verdict and a perfect code rendering are not the same achievement.

\textbf{Internal validity.} \textsc{OmniDiff} is timed in-process while every external tool is
timed as a whole process, so it pays no process-startup cost that the others pay in full. We state
this rather than correct for it, and report GumTree and BDiff both cold and warm so that a reader
can see how large the effect is for the two tools where it dominates. Per-tool mean runtimes are
additionally distorted by run order, whichever tool runs first per change absorbing cold-cache
cost; to minimise this effect we quote percentiles rather than means throughout, and every timing
number is the result of repeated runs.

The parts of this paper were not all measured at one state of the dataset and the tool. The tool
comparison, the cost comparison and \textsc{OmniDiff}'s line-level, speed and robustness figures
are over the \UnixDiffFixtures{} changes solved at an earlier stage and, for
\textsc{OmniDiff}, an earlier version; the node-level, painting-level and invariant figures are
over the \NumFixtures{} changes and the current version (footnote~\ref{fn:in-progress}). We keep
the two apart in the text rather than mix them in one table.

The ground truth is authored by the same people who built \textsc{OmniDiff}, which is a real risk
of bias in its favour. Three things limit it: the mappings are authored against the change itself
in a dedicated annotation tool with no tool output shown, the full annotation history is published
with the dataset so that any individual mapping can be inspected and disputed, and
Section~\ref{sec:oracle} compares our annotation against ground truth authored by other people,
for a different tool, before this one existed. That last check is the only one of the three that
does not depend on our own discipline, and it is the reason the dataset was extended with
Defects4J at all. The paintings and the content verdicts have no such check: they are made by the
same people, and we know of no external ground truth for how a code change should be painted or
for how much a change to a picture, a font or an archive changed.

The content timings are a single run per pair, where every code timing is repeated, so a content
percentile carries the noise of one measurement.

\textbf{External validity.} The fixtures are deliberately not a fully random sample of real
commits. They concentrate hard, structurally interesting changes, because easy ones exercise
nothing worth measuring, so every tail figure in this paper is a stress reading rather than an
estimate of what a developer waits for on a typical commit. Every accuracy result is also
reported per dataset (Figure~\ref{fig:agreement-by-dataset},
Table~\ref{tab:omnidiff-by-dataset}), so a reader can separate the two random samples, the
stratified one, and the Defects4J bug fixes from each other. The Defects4J units solved so far
are how far annotation has reached, not a draw, and they run shorter than the ones still
unsolved.

The painted subset is not randomly drawn either: painting is manual, and which changes have been
painted is an artifact of the order the annotators worked in. It runs to shorter changes than the
dataset as a whole, so the rates of RA2.1 and RA3.2 may move in either direction as painting
reaches larger changes. A living dataset also means that the figures of this paper describe one
day of it; the scripts that produce them are published with it and re-run on any later state.

The project's popularity says nothing about the quality of its ground truth, and its growth from
user reports is so far one content type (Section~\ref{sec:growth}).

Finally, the largest fixture holds \RobustnessMaxNodes{} AST nodes, well short of the Robust
target's \AstMax{}. The target is exercised only by the whole-corpus run of
Section~\ref{sec:omnidiff-evaluation}, whose budget and memory cap were applied with eight
processes sharing the machine, so a pair near either limit may have crossed it for want of a core
rather than of an algorithm. Re-running the timeouts on a quiet machine moved two of them into the
completed set and no more, and the memory kills recur across every commit of the same files, so
neither limit looks like an artefact of the sharing.

The judged content changes are fewer and less evenly spread than the code ones. They are the
sampled pairs annotation has reached so far, not a draw from them, and the families differ widely
in size: \ContentAccPicturesJudged{} pictures but \ContentAccPdfsJudged{} PDFs. We measured no
established tool on content (RA5.1), so the content results of Section~\ref{sec:omnidiff} have
no counterpart among the established tools.

\section{Related Work}
\label{sec:related}

\textbf{Deltas and their quality.} Section~\ref{sec:deltas} reviews the work on edit scripts.
What it shares is its criterion: a script is better when it is shorter, whether through block
moves~\cite{tichy-1984}, through more detected moves~\cite{mtdiff-2016}, or through more
compact tree scripts~\cite{ijm-2018,matsumoto-2019,falleri-2024}. This paper does not propose a better delta.
It measures, on changes with a human answer, how far the best delta is from what a reader should
be shown (RA2.1--RA2.4), and it finds that the distance is not closed by a better mapping alone.

\textbf{Visualisations and their quality.} The work closest in intent is the work that judges a
diff by its reader. srcDiff~\cite{decker-2020} argues that a delta should follow how developers
edit rather than minimise cost and supports this with a user study; ClDiff~\cite{cldiff-2018},
LSdiff~\cite{kim-notkin-2009} and the refactoring-aware review
tools~\cite{ge-2017,brito-2021,refactoringminer-2022} summarise a change at a higher level than
its edit script. We add to this line a statement of what a visualisation has to satisfy that is
independent of any tool (the invariants), a ground truth of the visualisation itself (the
paintings), and the observation that the same distinction governs content that is not code,
where perceptual image comparison~\cite{ssim-2004,yee-2004} has made it for pictures already.

\textbf{Tree-diffing tools.} Section~\ref{sec:background-tools} describes what each of the tools
compared here does; what matters for positioning this work is where each one's goal differs from
ours. GumTree~\cite{gumtreediff-2014} remains the reference point, and \textsc{OmniDiff}'s
move-recovery phase adapts its recovery-mappings heuristic directly: the difference is that
GumTree targets research use, with broad language and backend flexibility that stays valuable
well beyond what RA5.2's single dataset measures, while \textsc{OmniDiff} narrows to a
tree-sitter-only frontend and spends the saved generality on production targets.
difftastic~\cite{difftastic} and diffsitter~\cite{diffsitter} share that parsing layer but
optimise for a diff a human reads comfortably rather than for fidelity of the underlying node
mapping, so RA5.2 measures them against an objective they did not choose. Against all of them,
Unix \texttt{diff} and \texttt{git diff}~\cite{myers-1986} remain the baseline: on our dataset
they are the fastest measured, and at the line level they render \UnixDiffPerfectPct{}\% of the
changes perfectly, against \DifftasticPerfectPct{}\% for difftastic. BDiff's~\cite{bdiff} block
awareness does not move that metric, and neither does the choice among \texttt{git}'s four
algorithms~\cite{nugroho-histogram}.

\textbf{Ground truth and its ambiguity.} GumTree's own evaluation, and most that followed it,
score a tool against a single reference mapping per change, and the oracle of Alikhanifard and
Tsantalis~\cite{astdiffbenchmark} is the most carefully built of these. RA2.3, RA3.1 and RA3.2 are
an argument that this default understates what a ground truth can be asked: on the changes where
the correspondence is not one-to-one, where several mappings are equally correct, or where several
renderings are, a fixed reference records a correct answer as wrong. The dataset here departs from
the default in three ways---a multi-map group, which records an $N{:}M$ site and admits any
consistent one-to-one pairing of it; paintings in two presets, with the single painting as the
recorded finding that a change has one rendering; and invariants that hold all of it to one
convention---and all three were added because annotating real changes forced them, not by design.
We are not aware of a code-diffing ground truth that records any of the three.

\textbf{Empirical studies.} Dyer et al.~\cite{ast-api-1} and L\"{a}mmel et al.~\cite{ast-api-2}
mine AST nodes at scale to study API usage; Arafat and
Riehle~\cite{commit-size-distribution-paper} measure the commit-size distribution of open-source
software, the distribution \textsc{OmniDiff}'s Fast target is stated against; and
Defects4J~\cite{defects4j} collects real, reproducible bug fixes from open-source Java projects,
each a before/after file pair, and has since served as the file-pair corpus behind most
evaluations of AST differencing on Java. This paper's empirical study differs from these in
purpose rather than method: it measures what commits change, by content, and how large the
changes are, in order to set and then evaluate the targets of a diffing tool.

\textbf{Diffing content other than code.} \textsc{OmniDiff}'s container engine takes the model of
diffoscope~\cite{diffoscope}, a file as members, each compared by what it is, and applies it to a
different end. diffoscope reports every byte that differs, because for a reproducible
build~\cite{lamb-reproducible-2022} every byte matters; \textsc{OmniDiff} decodes only the members
whose hash changed and names each change by how much of it a reader would see, so a re-compressed
archive or a re-hinted font shows as unchanged. Git's \texttt{textconv}~\cite{gitattributes}
reaches the same formats through a converter per format that a user configures by path pattern;
\textsc{OmniDiff} recognises the format by its bytes, with no configuration, and diffs a picture
as pixels rather than as text. Its picture engine uses pixelmatch's perceptual
test~\cite{pixelmatch} per pixel and adds what a diff needs beyond a count of changed pixels:
where the changed regions are, and how the frames of an animation align. We are not aware of a
ground truth for content diffs that records how much a change changed; the judged changes
published with this paper are one.

\section{Conclusions}
\label{sec:conclusion}

This paper separated the delta of a change, which a machine executes, from its visualisation,
which a human reads, and measured the distance between the two on a living ground-truth dataset
of changes to code and to six families of other content. The distance is large for every content
type we measured. Git shows \ChangeCensusBinaryPct{}\% of the files commits change only as
differing (RA1.1), and \ContentVerdictNothingPct{}\% of judged changes to such content have a
byte delta and nothing deliberate to show (RA2.2). For code, a line script painted as it stands
breaks an invariant of visualisation on \InvariantLineScriptFixturesPct{}\% of changes, and a
tree script with the human's own mapping on \InvariantTreeScriptFixturesPct{}\% (RA2.1); on
\AmbiguityListPct{}\% of changes the preferred correspondence cannot be expressed with the four
standard operations at all (RA2.3), and on \CostHumanHigherPct{}\% the cheapest mapping is not
the one humans prefer (RA2.4). Neither the mapping nor its rendering is unique (RA3.1, RA3.2).
Established tools meet developers' speed requirements, the line-based ones within tens of
milliseconds (RA5.3), and fall short of the human-preferred rendering: no established tool that
parses most code changes renders much more than three fifths of them perfectly (RA5.2).
\textsc{OmniDiff} diffs \ContentCensusShownPct{}\% of changed files as text or by their content,
which leaves \ContentCensusLeftPct{}\% that a reader sees only as differing bytes, at a speed
that meets its target for pictures, catalogs and cursors and not yet for archives, fonts and
PDFs.

Based on the results reported in this paper, we recommend the following:

\begin{itemize}
  \item Evaluate a diffing tool against what a human would mark, not against the length of its
    edit script: a cheapest script and a preferred account differ measurably.
  \item Treat the projection from a mapping to what is displayed as a component of its own, with
    stated rules; a correct mapping does not make a readable diff.
  \item State the conventions of a hand-made ground truth as machine-checked invariants, and run
    them on every answer; they find the ground truth's own errors as well as a tool's.
  \item Record more than one correct answer where there is more than one, in the mapping and in
    the rendering, and record explicitly where there is exactly one.
  \item Report accuracy on what reaches the reader: visible nodes and painted bytes for code, and
    how much of a change is visible for other content.
  \item Diff a file by what it holds. \ChangeCensusBinaryPct{}\% of changed files are binary
    to git, and for those files a notice that the bytes differ is the whole of today's standard
    diff.
\end{itemize}

\bibliographystyle{ACM-Reference-Format}
\bibliography{references}

\end{document}
